% \documentclass[12pt]{article}
\documentclass[12pt, draft]{article}
% 
\usepackage{setspace}
\usepackage[labelfont=bf]{caption}
% \usepackage[margin=1in]{geometry}
% \usepackage[a4paper, total={6in, 10in}]{geometry}
\usepackage[margin=1in]{geometry}
\usepackage{amsmath,amssymb,amsthm,bm,mathtools}
\usepackage{algorithm}
\usepackage{algpseudocode}
\usepackage{enumitem}
\usepackage{hyperref}

% \documentclass{article}
\usepackage{graphicx} % Required for inserting images
\usepackage{hyperref}
% \usepackage[a4paper, total={5in, 10in}]{geometry}
% \usepackage{xcolor}
% \usepackage{color,soul}
\usepackage{listings}
\usepackage{booktabs}
% \usepackage{natbib}

% \usepackage{graphicx}
\usepackage{xcolor} % Required for colored text

% Define custom colors to match the image's muted green and red
\definecolor{impGreen}{rgb}{0.0, 0.5, 0.0}
\definecolor{worsRed}{rgb}{0.7, 0.0, 0.0}

% Citations
% \usepackage{natbib}
% Use the biblatex package, specifying the backend (biber is recommended)
\usepackage[style=numeric, sorting=nyt]{biblatex}
% Import your bibliography file (omit the .bib extension here)
\addbibresource{references.bib}


% \documentclass{article}
% \usepackage{booktabs}
% \usepackage{geometry}
% \geometry{a4paper, margin=1in}

% Math
\usepackage{amsmath}
\usepackage{amssymb}

% Subplots
\usepackage{graphicx}
\usepackage{subcaption}

% colored area
\usepackage{tcolorbox}
\usepackage{xcolor} % Required for defining custom colors
\definecolor{lightgray}{gray}{0.9} % Define a light gray color

% Highlights
\usepackage{color,soul}

% TODOs
% \usepackage{todonotes}
\usepackage[textsize=tiny]{todonotes}
% Indicator
\usepackage{dsfont}

% Line numbers
\usepackage{lineno}
\renewcommand{\linenumberfont}{\normalfont\bfseries\tiny\color{gray}}
\linenumbers

% 
% \documentclass{book} % Use the book or report class


\hypersetup{colorlinks=true,linkcolor=blue,citecolor=blue,urlcolor=blue}


\title{\vspace{-5mm}Fairness in Machine Learning for Real Estate Valuation:\\ Reducing Regressivity in Mass Appraisal \vspace{0mm}}
% Simple Regressivity Control for Mass Appraisal Pipelines
\author{Nicolás Acevedo Villena\\
{\small Operations Research Center, MIT, Cambridge, MA 02139. \texttt{nacevedo@mit.edu}}}
\date{
% \small
% Operations Research Center
\vspace{-8mm}
}

% theorem envs
\theoremstyle{plain}
\newtheorem{theorem}{Theorem}
\newtheorem{proposition}{Proposition}
\theoremstyle{definition}
\newtheorem{definition}{Definition}
\newtheorem{assumption}{Assumption}
\newtheorem{remark}{Remark}

\newcommand{\R}{\mathbb{R}}
\newcommand{\E}{\mathbb{E}}
\newcommand{\prox}{\mathrm{prox}}
\newcommand{\argmin}{\mathop{\mathrm{arg\,min}}}
\newcommand{\norm}[1]{\left\lVert #1 \right\rVert}
\newcommand{\ip}[2]{\left\langle #1,#2 \right\rangle}
\newcommand{\1}{\mathbf{1}}
\newcommand{\Cov}{\text{Cov}}
\newcommand{\Var}{\text{Var}}


% Set commands
\newcommand{\Y}{\mathcal{Y}}
\newcommand{\X}{\mathcal{X}}
\newcommand{\V}{\mathcal{V}}
\newcommand{\U}{\mathcal{U}}
\newcommand{\G}{\mathcal{G}}
\newcommand{\D}{\mathcal{D}}
\newcommand{\N}{\mathcal{N}}
\newcommand{\F}{\mathcal{F}}
\newcommand{\A}{\mathcal{A}}
\newcommand{\J}{\mathcal{J}}
\newcommand{\I}{\mathcal{I}}
\newcommand{\M}{\mathcal{M}}
\newcommand{\C}{\mathcal{C}}
\renewcommand{\L}{\mathcal{L}}
\renewcommand{\H}{\mathcal{H}}
\renewcommand{\P}{\mathcal{P}}
\renewcommand{\R}{\mathbb{R}}

\newcommand{\diag}{\mathbf{\operatorname{diag}}}
% \newcommand{\R}{\mathbb{R}}
% \newcommand{\1}{\mathbf{1}}

\newcommand{\muBold}{\boldsymbol{\mu}}
\renewcommand{\t}{^{\top}}
\newcommand{\qt}[1]{``#1''}

\newcommand{\Expected}[1]{\mathbb{E}[#1]}
\newcommand{\Prob}[1]{\mathbb{P}(#1)}

\newcommand{\red}[1]{{\color{red} #1}}

% Math sets commands


\doublespacing
\begin{document}
\maketitle
% \tableofcontents

\begin{abstract}

Mass appraisal is a critical tool for local governments: it helps determine property taxes, one of the major sources of local public revenue in the United States. In recent years, assessors have increasingly adopted Machine Learning (ML) methods in the mass appraisal process. However, the ML-based property valuation literature has primarily focused on predictive accuracy, while giving less attention to equity concerns in assessment. One of the main equity concerns in mass appraisal is \textit{regressivity}, a form of vertical inequity, whereby lower-value homes tend to be over-assessed while higher-value homes tend to be under-assessed. We propose a simple assessor-compatible regressivity-aware modification of standard regression loss functions that can be incorporated into regression models commonly used in practice by assessors (e.g., linear regression \& gradient boosting). Rather than proposing a new valuation paradigm, the modification is designed as a targeted correction that can be integrated into existing assessor workflows with minimal disruption. Specifically, the modification consists of a covariance-based regularization term that softly penalizes the dependence between assessment ratios and property values. We empirically show that these soft-constrained models can substantially mitigate regressivity while incurring little to no loss in out-of-sample accuracy. Finally, we propose a multi-objective optimization framework for model selection that identifies Pareto-efficient tradeoffs among competing performance measures (e.g., accuracy \& regressivity) on out-of-sample data.

\end{abstract}

% \begin{abstract}


% {\color{gray}
% Mass appraisal is a critical tool for local governments: it helps determine property taxes, one of the major sources of local public revenue in the United States. In recent years, assessors have increasingly adopted Machine Learning (ML) methods in the mass appraisal process. However, the ML-based property valuation literature has primarily focused on predictive accuracy, while giving less attention to equity concerns in assessment. One of the main equity concerns in mass appraisal is \textit{regressivity}, a form of vertical inequity, whereby lower-value homes tend to be over-assessed while higher-value homes tend to be under-assessed. We propose a simple, principled regressivity-aware modification of standard regression loss functions that can be incorporated into regression models commonly used in practice by assessors (e.g., linear regression \& gradient boosting). Specifically, the modification consists of a covariance-based regularization term that softly penalizes the dependence between assessment ratios and property values. We empirically show that these soft-constrained models can substantially mitigate regressivity while incurring little to no loss in out-of-sample accuracy. Finally, we propose a multi-objective optimization framework for model selection that identifies Pareto-efficient tradeoffs among competing performance measures (e.g., accuracy \& regressivity) on out-of-sample data.
% }






% % \hrule

% % Mass appraisal is a critical issue for local governments: it helps determine property taxes, their largest source of revenue in the US. In recent years, assessors have shifted towards the use of Machine Learning (ML) in their mass appraisal process. However, the ML-based property valuation literature has mainly focused solely on obtaining the most accurate predictive model, while neglecting potential collateral risks; e.g., securing a fair mass appraisal process. One of the main concerns in terms of fairness in mass valuation is that assessments tend to be \textit{regressive} (low-value homes are over-predicted, while high-value homes are under-predicted). We propose a simple regressivity-aware loss modification for general regression problems: a covariance-based regularization term that can be added to softly control unfair behaviors in many loss objectives of regression models used by assessors nowadays (e.g., linear regression \& gradient boosting). We empirically show that these soft-constrained models mitigate regressivity up to desired levels, while sacrificing bare to none predictive power. Finally, we propose a multi-objective optimization methodology for model selection, that aims to identify the optimal tradeoffs between different performance measures (e.g., accuracy \& regressivity) in out-of-samples sets.


% % which allows us to not just constraint regressivity in-sample, but also generalize this behavior correction to out-of-sample sets.   


% % desired trade-offs between accuracy and fairness. Further, we propose a convex robust optimization as model selection step, which  can secure hard fairness constraints on out-of-sample sets (empirically prove?).
%     % \ldots
% % \par\vspace{\baselineskip}\noindent \textit{Keywords:} fairness, regressivity, mass appraisal, multi-objective criteria, robustness
%     % \textit{Keywords:} 
% \end{abstract}

% \tableofcontents

% \section{Introduction}

\section{Introduction}

Mass appraisal is a critical component of local public finance: it helps local governments determine property taxes, often the largest source of local public revenue in the United States \cite{Berry2021}. Mass appraisal refers to the process of estimating the value of a group of real estate properties, as of a given date, for tax purposes. The process is done using common data, standardized methods, and statistical testing \cite{IAAO2025Mass}. Its importance extends beyond property owners for two reasons. First, assessments help determine property tax liabilities, which in turn finance roads, schools, and many other public services.  Second, they affect property values and housing costs more broadly, with consequences that extend to much of the population, both directly and indirectly \cite{Berry2021, Rootetal2023}. Therefore, securing an \textit{accurate} and equitable (\textit{fair}) assessment process is essential for local public finance and social welfare.

In recent years, local government assessors' offices have increasingly adopted Machine Learning (ML)-based methods in mass appraisal, motivated by their strong predictive performance, efficiency, and scalability in large valuation settings \cite{IAAOAITaskForce2022,BidansetRakow2022,Rootetal2023}. Commonly studied model classes in the mass appraisal literature include linear regression-based methods and tree ensembles \cite{Rootetal2023,Ostrikova2024,Sevgen2024100111}. In this setting, however, predictive gains alone are not sufficient: methods must also remain explainable, operationally feasible, and compatible with assessor workflows. Naturally, the growing use of ML raises important equity concerns in assessment. Two of the main concerns are whether ML-based valuation can produce \textit{fair} predictions across: (i) properties with similar characteristics (\textit{horizontal equity}), and (ii) properties of different price level strata (\textit{vertical equity}) \cite{Quintos2014,10578d49-78f8-3ef9-8277-9b6add2bd071,cornia2006property}. Although both concepts have long been studied, even before the recent wave of ML adoption, vertical equity remains a persistent challenge in practice; assessments often exhibit \textit{regressivity} \cite{McMillen2020,IAAO2023VerticalReview,McMillenSingh2023}. Hence, this paper focuses primarily on reducing vertical equity problems, and in particular on reducing regressivity in ML-based valuation.

\textit{Regressivity} arises when lower-valued properties tend to be over-assessed, while higher-valued properties tend to be under-assessed. This pattern has long been recognized in the mass appraisal literature and can have substantial distributional consequences for taxpayers \cite{McMillen2020}. \textcite{amornsiripanitch2020why} argues that owners of
lower-valued properties might face effective tax rates up to 50\% higher than those paid by owners of higher-valued properties. Data limitations, model misspecification, and assessment practices are often cited as contributing factors \cite{Berry2021}. The adoption of ML-based assessment workflows creates both an opportunity and a challenge in this setting. On one hand, these workflows provide a mathematical object that can be explicitly modified to mitigate regressivity (see, e.g., \textcite{CandoganHanLu2023}). On the other hand, standard objectives in regression models are designed to optimize average predictive accuracy, which tends to pull extreme predictions toward conditional averages. Thus, if valuation features are incomplete or imperfectly measured in the mass appraisal context, predictive models can reinforce regressive assessment patterns \cite{amornsiripanitch2020why, CandoganHanLu2023}.

% by overvaluing lower-priced properties and undervaluing higher-priced ones 

To evaluate assessment equity, we follow the standard framework used in ratio studies by the International Association of Assessing Officers (IAAO) \cite{IAAO2025ExposureRatio}. Many of these measures are based on the \textit{assessment-to-sale ratio}, defined as $r := \hat{P}/P$, where $\hat{P}$ denotes the assessed value and $P$ the market value of a property. Under regressive assessment, these ratios tend to decrease with property value: lower-tier properties are more likely to have ratios above one, while higher-tier properties are more likely to have ratios below one. Figure~\ref{fig:regressivity_motivation} illustrates this phenomenon. Motivated by this pattern, we study an in-processing approach that penalizes systematic dependence between assessment ratios and property values, with the goal of reducing regressivity while preserving predictive accuracy.

\begin{figure}[!h]
    \centering 
    \includegraphics[width=0.7\linewidth]{Fairness Price/CCAO meeting/img/Baseline LightGBM_0_motivation.pdf}
    \caption{Illustration of regressive behavior in mass appraisal. Assessment-to-sale ratios $r$ tend to decrease with property value $P$: lower-value properties are more likely to be over-assessed ($r > 1$), while higher-value properties are more likely to be under-assessed ($r < 1$).}
    \label{fig:regressivity_motivation}
\end{figure}

More broadly, this paper helps bridge two bodies of work that have largely developed in parallel: the ML-based mass-appraisal literature, which has mainly emphasized predictive performance model comparison in property valuation \cite{Rootetal2023,Zilli2024}, and the fairness-in-ML literature, which has developed general frameworks for fair regression, fairness regularization, and fairness--accuracy tradeoff analysis in predictive models \cite{CatonHaas2024,AgarwalDudikWu2019,FitzsimonsEtAl2019,lee2022maximal,LiEtAl2022Fairness}. This connection is especially natural in mass appraisal because regressivity can be viewed as a systematic dependence between assessment ratios and property values, allowing an assessor-defined notion of vertical inequity \cite{IAAO2025ExposureRatio} to be addressed using tools from fairness-aware regression \cite{CandoganHanLu2023}.

\noindent \textbf{Contributions.} Our main contributions are as follows:
\begin{itemize}
    \item We develop a simple, in-processing approach for mitigating regressivity in ML-based mass appraisal by augmenting standard regression losses with a soft covariance-based regularization term involving transformations of the assessment-to-sales ratio $r$ and the property market value $P$.\footnote{Related dependence-regularization ideas appear in the fair-ML literature; see, e.g., \cite{lee2022maximal,LiEtAl2022Fairness,ScutariPaneroProissl2022}.}
    
    \item We show how this approach can be incorporated into regression models commonly used in practice by assessors, including linear regression and gradient boosting, thereby providing a practical way to reduce vertical inequity without requiring a complete redesign of existing valuation workflows.
    
    \item Using a real-world mass-appraisal pipeline from the Cook County Assessor's Office, we show empirically that the proposed corrections can improve equity--accuracy tradeoffs under standard assessment metrics, and in some settings improve both dimensions simultaneously.
    
    \item We formalize a multi-objective model-selection framework for calibration that targets Pareto-efficient tradeoffs between predictive accuracy and equity measures. In particular, the framework can recover solutions that prioritize predictive performance subject to equity requirements, and it also admits an extension based on convex combinations of candidate models to achieve improved tradeoffs. \footnote{For related Pareto/fairness-frontier formulations and model-stacking ideas, see, e.g., \cite{martinez2020minimax,little2022fairness}.}
\end{itemize}

The remainder of the paper reviews the relevant background on mass-appraisal equity and fairness-aware regression, develops the proposed regularization and model-selection framework, and evaluates the resulting methods on a real-world assessment pipeline.





% \section{Introduction}

% \begin{itemize}
%     % \item What is Mass Appraisal and why is it (so) relevant? (this part should be relatively short?) % DONE
%     % \item The change of paradigm to ML-based mass appraisal: ensemble models being in the top used with high accuracy (extensive reviews and papers, although not so good) % DONE
%     % \item High-accuracy models are good for exactly that: reducing error (and maybe dispersion), but they may increase vertical inequality significantly (regressive behavior mainly).
%     % \item The problem of regressivity, What it is exactly, its consequences \& the extensiveness of its literature. (not much for ML [huge gap here?], but still) % DONE
%     \item What are the current approaches to solve regressivity? Haihao, Quintos, \& some other maybe. What are its advantages \& limitations? Why do we need this approach? THIS IS FOR LITERATURE REVIEW
%     \item  What is our contribution to solve this problem? Its advantages/contributions? Limitations? -> Probably the generalization to any pipeline? Its simplicity, and direct correction power: achieves fairness with highly accurate models accuracy. Can we do it for RF? (THIS IS FOR OTHER SECTIONS?)
%     % \item Do we introduce a robust model selection framework here? If we add it in contributions it need to be motivated? I don't think so. Its sufficient to mention.
%     \item Structure of the paper?
    
% \end{itemize}

% \hrule 
% \newpage

% % \section{Introduction}

% % v2
% Mass appraisal is a critical component of local public finance: it helps local governments determine property taxes, often the largest source of local public revenue in the United States \cite{Berry2021}. Mass appraisal refers to the process of estimating the value of a group of real estate properties, as of a given date, for tax purposes. The process is done using common data, standardized methods, and statistical testing \cite{IAAO2025Mass}. Its importance extends beyond property owners for two reasons. First, assessments help determine property tax liabilities, which in turn finance roads, schools, and many other public services.  Second, they affect property values and housing costs more broadly, with consequences that extend to much of the population, both directly and indirectly \cite{Berry2021, Rootetal2023}. Therefore, securing an \textit{accurate} and equitable (\textit{fair}) assessment process is essential for local public finance and social welfare.

% % V1
% % Mass appraisal is a critical component of local public finance for local governments, especially municipalities \cite{Rootetal2023}. It refers to the set of methods used by local government offices to determine the value of large numbers of real estate property. It is a crucial step in determining property taxes---the largest source of revenue for local governments in the USA---used to fund roads, schools, and many other public services \cite{Berry2021}. Consequently, property taxes affect not only real estate owners, but also much of the population, either directly or indirectly, through property prices, rents, and public services. Therefore, securing an accurate and fair assessment process is essential for the local economy and social welfare.

% % V2
% In recent years, local governments' assessors offices have increasingly adopted Machine Learning (ML)-based methods in mass appraisal, motivated by their strong predictive performance, efficiency, and scalability in large valuation settings \cite{IAAOAITaskForce2022,BidansetRakow2022,Rootetal2023}. Commonly studied model classes in the mass appraisal literature include linear regression-based methods and tree ensembles \cite{Rootetal2023,Ostrikova2024,Sevgen2024100111}. Naturally, the growing use of ML raises important equity concerns in assessment. Two of the main concerns being whether ML-based valuation can produce \textit{fair} predictions across: (i) properties with similar characteristics (\textit{horizontal equity}), and (ii) properties of different price level strata (\textit{vertical equity}) \cite{Quintos2014,10578d49-78f8-3ef9-8277-9b6add2bd071,cornia2006property}. Although both concepts have long been studied, even before the recent wave of ML adoption, vertical equity remains a persistent challenge in practice; assessments often exhibit \textit{regressivity} \cite{McMillen2020,IAAO2023VerticalReview,McMillenSingh2023}. This paper focuses primarily on vertical equity, and in particular on reducing the problem of regressivity in ML-based valuation.

% % V1
% % In recent decades, local governments assessors offices have shifted toward the use of Machine Learning (ML) models, particularly linear regression-based models and tree ensembles, to assist in the mass appraisal process, given their higher predictive performance compared to individual property assessments \cite{Rootetal2023, Ostrikova2024,Sevgen2024100111}. Naturally, this raises concerns, with two of the major ones being the capacity of ML to produce \textit{fair} predictions across: (i) properties with similar characteristics (\textit{horizontal equity}), and (ii) properties of different price levels (\textit{vertical equity}) \cite{Quintos2014}. These two concepts have been extensively studied in the literature \cite{10578d49-78f8-3ef9-8277-9b6add2bd071,cornia2006property}, even before the ML era, but in practice vertical equity remains a challenge; assessments tend to be \textit{regressive} \cite{McMillen2020}. 

% % V2
% \textit{Regressivity} arises when lower-valued properties tend to be over-assessed, while higher-valued properties tend to be under-assessed. This pattern has long been recognized in the mass appraisal literature and can have substantial distributional consequences for taxpayers \cite{McMillen2020}. \textcite{amornsiripanitch2020why} argues that owners of
% lower-valued properties might face effective tax rates up to 50\% higher than those paid by owners of higher-valued properties. Data limitations, model misspecification, and assessment practices are often cited as contributing factors \cite{Berry2021}. The adoption of ML-based assessment workflows creates both an opportunity and a challenge in this setting. On one hand, these workflows provide a mathematical object that can be explicitly modified to mitigate regressivity (see, e.g., \textcite{CandoganHanLu2023}). On the other hand, standard objectives in regression models are designed to optimize average predictive accuracy, which tends to pull extreme predictions toward conditional averages. In the mass appraisal context, when valuation features are incomplete or imperfectly measured, this can reinforce regressive assessment patterns by overvaluing lower-priced properties and undervaluing higher-priced ones \cite{amornsiripanitch2020why, CandoganHanLu2023}.

% % Under standard squared-error objectives, regression models target conditional averages. In mass appraisal, when valuation features are incomplete or imperfectly measured, this can shrink predictions toward local averages and contribute to overvaluation at the low end of the market and undervaluation at the high end .

% % The increasing use of ML creates both a challenge and an opportunity in this setting: more flexible models can improve predictive performance and horizontal equity, but they also make it especially important to explicitly control systematic value-related distortions in assessments rather than treating them only as ex-post diagnostic outcomes.

% % When trained under standard squared-error objectives, regression models target conditional averages; when valuation features are incomplete or imperfectly measured, this can shrink predictions toward local averages and, in mass appraisal, contribute to overvaluation at the low end of the market and undervaluation at the high end \cite{CandoganHanLu2023,amornsiripanitch2020why}.



% % % V1
% % \textit{Regressivity} occurs when properties at the lower-end of the market tend to be over-assessed (assessed price > market price), while properties at the higher-end of the market tend to be under-assessed (assessed price < market price). This regressive pattern in assessments is argued to be prevalent throughout the United States. Its impact can be substantial; owners of
% % lower-valued properties might face effective tax rates up to 50\% higher than those paid by owners of
% % higher-valued properties \cite{amornsiripanitch2020why}. The limitations of the data and assessment methods used are often cited as reasons for these inequities \cite{Berry2021}. The adoption of ML-based assessment workflows opens an opportunity to correct these regressive behaviors mathematically (see, e.g., \textcite{CandoganHanLu2023}), but at the same time it can further enforce regressive patters in regression settings (\todo{CITE}).


% To evaluate assessment equity, we follow the standard framework used in ratio studies by the International Association of Assessing Officers (IAAO) \cite{IAAO2025ExposureRatio}. Many of these measures are based on the \textit{assessment-to-sale ratio}, defined as $r := \hat{P}/P$, where $\hat{P}$ denotes the assessed value and $P$ the market value of a property. Under regressive assessment, these ratios tend to decrease with property value: lower-tier properties are more likely to have ratios above one, while higher-tier properties are more likely to have ratios below one. Figure~\ref{fig:regressivity_motivation} illustrates this phenomenon. Motivated by this pattern, we study an in-processing approach that penalizes systematic dependence between assessment ratios and property values, with the goal of reducing regressivity while preserving predictive accuracy.

% \begin{figure}[!h]
%     \centering 
%     \includegraphics[width=0.7\linewidth]{Fairness Price/CCAO meeting/img/Baseline LightGBM_0_motivation.pdf}
%     \caption{Illustration of regressive behavior in mass appraisal. Assessment-to-sale ratios $r$ tend to decrease with property value $P$: lower-value properties are more likely to be over-assessed ($r > 1$), while higher-value properties are more likely to be under-assessed ($r < 1$).}
%     \label{fig:regressivity_motivation}
% \end{figure}


% More broadly, this paper helps bridge two bodies of work that have largely developed in parallel: the ML-based mass-appraisal literature, which has mainly emphasized predictive performance and model comparison in property valuation \cite{Rootetal2023,Zilli2024}, and the fairness-in-ML literature, which has developed general frameworks for fair regression, fairness regularization, and fairness--accuracy tradeoff analysis in predictive models \cite{CatonHaas2024,AgarwalDudikWu2019,FitzsimonsEtAl2019,lee2022maximal,LiEtAl2022Fairness}. This connection is especially natural in mass appraisal because regressivity can be viewed as a systematic dependence between assessment ratios and property values, allowing an assessor-defined notion of vertical inequity \cite{IAAO2025ExposureRatio} to be addressed using tools from fairness-aware regression \cite{CandoganHanLu2023}.

% % % V1
% % More broadly, this work aims to help bridge two literatures that have developed largely in parallel: the ML-based mass-appraisal literature, which has mainly emphasized predictive performance in this application, and the fairness-in-ML literature, which has developed general tools for controlling undesirable dependence structures in predictive models. In the present setting, this connection is especially natural because regressivity can be interpreted as a systematic dependence between assessment ratios and property values.

% % \medskip

% \noindent \textbf{Contributions.} Our main contributions are as follows:
% \begin{itemize}
%     % \item We develop, for ML-based mass appraisal, a simple in-processing approach to mitigating regressivity by augmenting standard regression losses with a soft covariance-based regularization term involving transformations of the assessment-to-sales ratio $r$ and the property market value $P$.

%     \item We develop a simple, in-processing approach for mitigating regressivity in ML-based mass appraisal by augmenting standard regression losses with a soft covariance-based regularization term involving transformations of the assessment-to-sales ratio $r$ and the property market value $P$.\footnote{Related dependence-regularization ideas appear in the fair-ML literature; see, e.g., \cite{lee2022maximal,LiEtAl2022Fairness,ScutariPaneroProissl2022}.}
    
    
%     % \item We derive concrete instantiations of this approach for regression models commonly used in practice by assessors, including linear regression and gradient boosting, thereby providing a practical way to reduce vertical inequity without requiring a complete redesign of existing valuation workflows.

        
%     \item We show how this approach can be incorporated into regression models commonly used in practice by assessors, including linear regression and gradient boosting, thereby providing a practical way to reduce vertical inequity without requiring a complete redesign of existing valuation workflows.

        
%     \item We show how this approach can be incorporated into regression models commonly used in practice by assessors, including linear regression and gradient boosting, thereby providing a practical way to reduce vertical inequity without requiring a complete redesign of existing valuation workflows.
    
%     \item Using a real-world mass-appraisal pipeline from the Cook County Assessor's Office, we show empirically that the proposed corrections can improve equity--accuracy tradeoffs under standard assessment metrics, and in some settings improve both dimensions simultaneously.

%     % \item Using a real-world mass-appraisal pipeline from the Cook County Assessor's Office, we empirically demonstrate that the proposed corrections can substantially improve equity--accuracy tradeoffs under standard assessment metrics, and in some cases improve both dimensions simultaneously.
    
%     % \item We formalize a multi-objective framework for model calibration and selection that targets Pareto-efficient tradeoffs between predictive accuracy and equity measures. We also show how this framework can be extended through convex combinations of candidate models to obtain improved tradeoffs.\footnote{For related Pareto/fairness-frontier formulations and model-stacking ideas, see, e.g., \cite{martinez2020minimax,little2022fairness}.}

%     \item We formalize a multi-objective model-selection framework for calibration that targets Pareto-efficient tradeoffs between predictive accuracy and equity measures. In particular, the framework can recover solutions that prioritize predictive performance subject to equity requirements, and it also admits an extension based on convex combinations of candidate models to achieve improved tradeoffs. \footnote{For related Pareto/fairness-frontier formulations and model-stacking ideas, see, e.g., \cite{martinez2020minimax,little2022fairness}.}
% \end{itemize}

% % OLD CONTRIBUTIONS
% % \medskip
% % \noindent \textbf{Contributions.} Our main contributions are as follows:
% % \begin{itemize}
% %     \item We introduce a simple, in-processing approach for mitigating regressivity in ML-based mass appraisal by augmenting standard regression losses with a soft covariance-based regularization term involving transformations of the assessment-to-sales ratio $r$ and the property market value $P$.
    
% %     \item We show how this approach can be incorporated into regression models commonly used in practice by assessors, including linear regression and gradient boosting, thereby providing a practical way to reduce vertical inequity without requiring a complete redesign of existing valuation workflows.
    
% %     \item Using a real-world mass-appraisal pipeline from the Cook County Assessor's Office, we empirically demonstrate that the proposed corrections can substantially improve equity--accuracy tradeoffs under standard assessment metrics, and in some cases improve both dimensions simultaneously.
    
% %     \item We formalize a multi-objective model-selection framework for calibration that targets Pareto-efficient tradeoffs between predictive accuracy and equity measures. In particular, the framework can recover solutions that prioritize predictive performance subject to equity requirements, and it also admits an extension based on convex combinations of candidate models to achieve improved tradeoffs.
% % \end{itemize}

% The remainder of the paper reviews the relevant background on mass-appraisal equity and fairness-aware regression, develops the proposed regularization and model-selection framework, and evaluates the resulting methods on a real-world assessment pipeline.

% % \hrule


% % OLD INTRODUCTION SECTION

% % Mass appraisal is a critical component of local public finance for local governments, especially municipalities \cite{Rootetal2023}. It refers to the set of methods used by local government offices to determine the value of large numbers of real estate property. It is a crucial step in determining property taxes---the largest source of revenue for local governments in the USA---used to fund roads, schools, and many other public services \cite{Berry2021}. Consequently, property taxes affect not only real estate owners, but also much of the population, either directly or indirectly, through property prices, rents, and public services. Therefore, securing an accurate and fair assessment process is essential for the local economy and social welfare.


% % In recent decades, local governments assessors offices have shifted toward the use of Machine Learning (ML) models, particularly linear regression-based models and tree ensembles, to assist in the mass appraisal process, given their higher predictive performance compared to individual property assessments \cite{Rootetal2023, Ostrikova2024,Sevgen2024100111}. Naturally, this raises concerns, with two of the major ones being the capacity of ML to produce \textit{fair} predictions across: (i) properties with similar characteristics (\textit{horizontal equity}), and (ii) properties of different price levels (\textit{vertical equity}) \cite{Quintos2014}. These two concepts have been extensively studied in the literature \cite{10578d49-78f8-3ef9-8277-9b6add2bd071,cornia2006property}, even before the ML era, but in practice vertical equity remains a challenge; assessments tend to be \textit{regressive} \cite{McMillen2020}. 


% % \textit{Regressivity} occurs when properties at the lower-end of the market tend to be over-assessed (assessed price > market price), while properties at the higher-end of the market tend to be under-assessed (assessed price < market price). This regressive pattern in assessments is argued to be prevalent throughout the United States. Its impact can be substantial; owners of
% % lower-valued properties might face effective tax rates up to 50\% higher than those paid by owners of
% % higher-valued properties \cite{amornsiripanitch2020why}. The limitations of the data and assessment methods used are often cited as reasons for these inequities \cite{Berry2021}. The adoption of ML-based assessment workflows opens an opportunity to correct these regressive behaviors mathematically (see, e.g., \textcite{CandoganHanLu2023}), but at the same time it can further enforce regressive patters in regression settings (\todo{CITE}).

% % To measure the equity level of a model, we follow the standard equity definitions as provided by the International Association of Assessing Officers (IAAO). In their Standard on Ratio Studies, the \textcite{IAAO2025ExposureRatio}\todo{check name and format} proposes a collection of metrics for both horizontal equity (e.g., the coefficient of dispersion) and vertical equity (e.g., the price-related differential), which are defined in (\todo{WHERE?}). The IAAO further defines the standard intervals within which the equity measures should lie \cite{iaao2025ratio}, which define our reference points. Most of these metrics rely on what is called the \textit{Assessment-to-Sales ratio}; that is, the ratio between $\Hat{P}$, the assessed value of a property, and $P$ its real market value. For simplicity, we denote this ratio value as \begin{equation}    
% % r(\hat{P}, P):= {\hat{P}}/{P}.
% % \end{equation}
% % This quantity, and the statistics build around it, help measure regressivity in different ways after a mass appraisal process. To illustrate this phenomenon, see the example in Figure \ref{fig:regressivity_motivation}. In this plot, there is a clear negative slope between some vector of ratios $r$ and its respective vector of market values $P$, which strongly suggest a regressive behavior on the assessments. Inspired by this illustration, our approach focuses on correcting the linear trends between $r$ and $P$, and some other relationship measures that can capture more complex relationships between them. Specifically, we take a covariance-based penalization approach between these two quantities.
% % \begin{figure}
% %     \centering 
% %     \includegraphics[width=0.7\linewidth]{Fairness Price/CCAO meeting/img/Baseline LightGBM_0_motivation.pdf}
% %     \caption{Regressive behavior between the log-market value $P$, and the Assessment-to-Sales ratios $r$. Lower-tier properties tend to be over-assessed ($r > 1$), while high-tier properties tend under-assessed ($r < 1$).}
% %     \label{fig:regressivity_motivation_2}
% % \end{figure}

% % \textbf{Contributions.} We make the following contributions: (i) To the best of our knowledge, this work is the first step in bridging the gap between the ML-based mass appraisal literature (usually focus on finding the most accurate models for this specific application), and the literature on fairness in machine learning (usually more theory-heavy), (ii) We propose an in-processing approach to help mitigate regressive trends in ML-based mass appraisal by adding a soft covariance-based constraint modification between transformations of $r$ and $y$ to general regression losses used by assessors (see \textcite{lee2022maximal} for related methodology). (iii) We empirically demonstrate, using a real-world mass appraisal pipeline (Cook County Assessors Office), that these corrections can achieve the desired accuracy-fairness tradeoffs—and, in some cases, improve both—under standard measures, using Pareto-based criteria, even as a post-processing step (see, e.g., \textcite{martinez2020minimax, little2022fairness}). (iv) Finally, we formalize a multi-objective model selection criteria that aims to produce a specific type of Pareto-optimal solution during the calibration stage: best accuracy subject to equity constraints, which can be done either by directly sorting the evaluated models and then selecting the best one, or it can be done by extending the selection process to a stacking of models, which allows a convex optimization formulation for many of the measures, which can achieve even better tradeoffs between accuracy and equity than single-model selection (see \textcite{lee2022maximal, little2022fairness} for related post-processing corrections). 

% % \textbf{}

% % MODEL SELECTION

% % that our model modification can also be used, as a simple post-processing step, to correct regressive patterns in models originally trained without such corrections, achieving performance similar to that of the original models while significantly improving standard equity measures.  







% % \textbf{Approach. } We propose a soft covariance-constrained approach to mitigate regressive patterns in general ML settings (see \textcite{lee2022maximal} for related methodology). Specifically, given a loss function $\L_n$ (e.g., mean squared error), we add a linear relationship regularization term between the ratio vector $r \in \R^n$ and market price vector $y\in \R^n_{++}$ in some learning problem, i.e., \vspace{-2mm}
% % \begin{equation}
% % \label{opt:unconstrained_learner}
% % \min_{f\in \H} \L_n(f) + \rho\frac{1}{2}{\textbf{Cov}}(r(f), y)^2, 
% % % := \frac{1}{2n}\sum_{i=1}^n (f(x_i) - z_i)^2 + \frac{\rho}{2}\left(\frac{1}{n}\sum_{i=1}^n(e_i - \Bar{e})(y_i - \Bar{y})\right)^2. %\qquad\text{ 
% % \vspace{-1mm}
% % \end{equation}
% % where $f$ is some mapping in a hypothesis class $\H$ (e.g., linear functions or tree ensembles) that defines predicted prices $\hat{y}=f(X)$, where $X \in \R^{n\times p}$ is a matrix of $p$ features (e.g., number of bedrooms, access to transport, etc.) for $n$ different properties. After{\color{white}-}training,{\color{white}-}we aim to robustly select a well-performing weight $\rho>0$, with respect to accuracy and fairness requirements, by defining an appropriate Pareto-based selection criterion (see, e.g., \textcite{martinez2020minimax, little2022fairness}). 


% % % Let $y \in \R^n_{++}$ denote the vector of market values for these $n$ properties. Let us further define the vector of log-prices $z \in \R^n$, and the vector of log-price predictions $\hat{z} \in \R^n$, obtained via a learned mapping $f:\R^p \rightarrow \R$ applied row-wise, and the estimation of market values $\hat{y}$ through exponentiation of $\hat{z}$. Finally, we define the \emph{log residuals} as \vspace{-2mm}
% % % $$
% % % e := \hat z - z \quad \Longrightarrow \quad e_i = \log(r_i), \quad \forall i \in [n]. \vspace{-2mm}
% % % $$
% % % Let $\mathcal H$ be a predefined hypothesis class. Then (in principle) we aim to learn an optimal correlation-aware (see \cite{cho2021maximal}) regression model via a linear dependency correction term $C_n(f)^2$:
% % % \begin{equation}
% % % \label{opt:unconstrained_learner_old}
% % % \min_{f\in \H} \L_n(f) + \rho\frac{1}{2}C_n(f)^2  := \frac{1}{2n}\sum_{i=1}^n (f(x_i) - z_i)^2 + \frac{\rho}{2}\left(\frac{1}{n}\sum_{i=1}^n(e_i - \Bar{e})(y_i - \Bar{y})\right)^2. %\qquad\text{ 
% % % \end{equation}
% % We further propose different proxies for the empirical covariance term $C_n(f)$, including a separable upper bound for it. The goal of this soft-constrained problem is to reduce the regressive behavior of models while using the same function class $\H$. Thus, the penalization parameter $\rho >0$ has to be selected in an independent process by maximizing accuracy on an out-of-sample set, similar to the FairStack pipeline in \cite{little2022fairness}.

% % Note that since the calibration of the weight parameter can be done in an independent step of the assessors' mass appraisal workflow, it does not add complexity to the model selection step. The selection of the new parameter $\rho$ can be done in a robust way in this special case of time-dependent data, by choosing the $\rho >0$ that produces the desired tradeoffs between the accuracy and equity metrics (see \cite{martinez2020minimax}) in some aggregated sense of the performance on validation sets constructed via rolling-origin forecasting resampling \cite{tashman2000out}.

% % \newpage

% % Assessment equity is a primary concern for municipalities for taxes purposes; the major concerns being a \textit{fair assessment} level across: (a) properties of different price-levels (\textit{vertical equity}), and (b) properties with similar characteristics (\textit{horizontal equity}) \cite{quintos2014improving}. The main caveat for accurately identifying the heterogeneity of the parcels in certain region is that the unsold properties need to be assessed (unknown market prices), and these estimations need to maintain consistency; for example, similar properties in the same block/area should be assigned similar prices. Given the sensitiveness of the matter, we constraint ourselves to the standard measures for equity defined by the International Association of Assessing Officers (IAAO); who proposes a collection of metrics for both horizontal and vertical equity \cite{iaao2025ratio} (e.g. the coefficient of dispersion and the price-related differential). Most of these metrics rely on what they call the \textit{appraisal ratio}, defined in terms of the prices and its predictions \todo{I am missing something about consequences and impact of taxes here, or at least some references.}
% % $$r(PV, MV):= \frac{PV}{MV}, \qquad {MV}:= \text{market value}, \qquad PV:= \text{predicted value. }$$
% % IAAO also defines the standard for the intervals where the equity measures should lie on, which we treat as hard constraints. In general, if for the results of a predictive model the majority of low-valued properties satisfy $r > 1$ (over-predicted), while the majority of the high-values properties satisfy $r < 1$ (under-predicted), we would say that our predictive model exhibits \textit{regressivity}. In the opposite case, we will say that it exhibits \textit{progressivity}. 

% % \textbf{Case of Study.} We focus on the regressive behavior exhibited by the predictive models used by the Cook County Assessor's Office (CCAO). The CCAO has an extensive open-access repository \cite{ccao_data_2026} with the datasets, code, analysis, and guides through the Mass Appraisal process they use for predicting property values in the Cook County, Chicago, IL. They used linear models for prediction between the years 2009 and 2018, and transitioned towards Gradient Boosting methods since 2019. Currently, the CCAO uses the state-of-the-art Newton-like gradient boosting method LightGBM \cite{lightgbm}.

% % \textbf{Approach. } Let us denote $X \in \R^{n\times p}$ denote a matrix of  $p$ demographic features for $n$ different properties (e.g., number of bedrooms, accessibility to transport, etc.) Let $y \in \R^n_{++}$ denote the vector of market values for each of these $n$ properties. Let us further define the vector of log-prices $z \in \R^n_{++}$, and  log-prices predictions vector $\hat{z}$, obtained via some learned mapping $f:\R^n \rightarrow \R$, and the estimation of market values $\hat{y}$ through exponentiation of $\hat{z}$. Finally, we define the \emph{log residuals} of the estimation of log-prices is defined as
% % $$
% % e := \hat z - z = f(x) - z  \Longrightarrow e_i = \log(r_i), \quad \forall i \in [n].
% % $$
% % Let $\mathcal H$ be a predefined hypothesis class, then (in principle) we aim to learn the optimal vertical-equity aware regression problem, via linear dependency correction $C_n(f)^2$ term
% % \begin{equation}
% % \label{opt:unconstrained_learner}
% % \min_{f\in \H} \L_n(f) + \rho\frac{1}{2}C_n(f)^2  := \frac{1}{2n}\sum_{i=1}^n (f(x_i) - z_i)^2 + \frac{\rho}{2}\left(\frac{1}{n}\sum_{i=1}^n(e_i - \Bar{e})(y_i - \Bar{y})\right)^2. %\qquad\text{ 
% % \end{equation}
% % We further propose different proxy's for the empirical covariance term $C_n(f)$, including a separable upper bound for it. The goal of this soft-constrained problem is to reduce the regressive behavior of models, using their same function class $\H$. Thus, the selection of the penalization parameters $\rho >0$ has to be selected in an independent process as the maximization of accuracy in out-of-sample set.

% % Note that, since the calibration of the weight parameter can be done in an independent step of assessors' Mass Appraisal workflow, it does not add more complexity to the model selection step. The new parameter $\rho$ selection can be done in a robust way in this special case of time-dependent data, by selecting the $\rho >0$ that produces the desired tradeoffs between the accuracy and equity metrics in some aggregated sense of the performance on validation sets constructed via rolling-origin forecasting resampling.




\section{Related Work}

\subsection{ML in Mass Appraisal and Assessor Workflows}

Recent reviews show that mass appraisal and real-estate valuation research have moved steadily toward automated valuation models (AVMs) and machine learning (ML) methods, with multiple regression, tree ensembles, and neural-network-based models among the most frequently studied approaches \cite{WangLi2019,Rootetal2023,Zilli2024,purdueAssessingTechniques}. Across this literature, model evaluation is still dominated by predictive criteria---typically accuracy and model comparison across algorithms---rather than by equity as an explicit model-design target \cite{WangLi2019,Rootetal2023,Zilli2024}. In particular, recent reviews identify Random Forest, Gradient Boosting, and related ensemble methods as some of the most commonly used and best-performing model classes in modern mass appraisal studies \cite{Rootetal2023,Sevgen2024100111,Zilli2024}. This accuracy-centered orientation leaves open a central question for the present paper: how should modern valuation models be designed and selected when equity is also a primary objective?

From an assessor perspective, however, adoption depends on more than predictive fit alone. The assessor-oriented AI literature emphasizes explainability, defensibility, administrative feasibility, and integration into office workflows as central conditions for real-world use \cite{IAAOAITaskForce2022}. Consistent with this, \textcite{BidansetRakow2022} report that AVM adoption in assessment offices remains partial and uneven, with perceived barriers including communication challenges, technical capacity, and budget constraints. As a result, the relevant literature for the present paper is not only the literature on which models predict most accurately, but also the literature on how advanced models can be incorporated into assessor practice in a way that remains operationally useful and publicly defensible \cite{IAAOAITaskForce2022,BidansetRakow2022}. These practical constraints make low-disruption mitigation strategies especially relevant in mass appraisal: an equity intervention is more likely to matter in practice when it preserves familiar model classes and fits within existing office workflows.

% \subsection{ML in Mass Appraisal and Assessor Workflows}

% Recent reviews show that mass appraisal and real-estate valuation research has moved steadily toward automated valuation models (AVMs) and machine learning (ML) methods, with multiple regression, tree ensembles, and neural-network-based models among the most frequently studied approaches \cite{WangLi2019,Rootetal2023,Zilli2024,purdueAssessingTechniques}. Across this literature, model evaluation is still dominated by predictive criteria---typically accuracy and model comparison across algorithms---rather than by equity as an explicit model-design target \cite{WangLi2019,Rootetal2023,Zilli2024}. In particular, recent reviews identify Random Forest, Gradient Boosting, and related ensemble methods as some of the most commonly used and best-performing model classes in modern mass appraisal studies \cite{Rootetal2023,Sevgen2024100111,Zilli2024}. This accuracy-centered orientation leaves open a central question for the present paper: how should modern valuation models be designed and selected when equity is also a primary objective?

% From the assessor perspective, however, adoption depends on more than predictive fit alone. The assessor-oriented AI literature emphasizes explainability, defensibility, administrative feasibility, and integration into office workflows as central conditions for real-world use \cite{IAAOAITaskForce2022}. Consistent with this, \textcite{BidansetRakow2022} report that AVM adoption in assessment offices remains partial and uneven, with perceived barriers including communication challenges, technical capacity, and budget constraints. As a result, the relevant literature for the present paper is not only the literature on which models predict best, but also the literature on how advanced models can be incorporated into assessor practice in a way that remains operationally useful and publicly defensible \cite{IAAOAITaskForce2022,BidansetRakow2022}.





\subsection{Equity in Mass Appraisal: Horizontal Equity, Vertical Equity, and Regressivity}

Equity has long been treated as a core objective of mass appraisal rather than as a secondary byproduct of predictive performance \cite{IAAO2013Ratio,IAAO2025Mass,Quintos2014}. In assessor practice, this objective is commonly decomposed into \emph{horizontal equity}, which concerns the uniform treatment of similarly situated properties, and \emph{vertical equity}, which concerns the absence of systematic price-related bias across the value distribution \cite{10578d49-78f8-3ef9-8277-9b6add2bd071,cornia2006property,McMillenSingh2023}. The basic diagnostic framework remains the sales-ratio study: current IAAO standards continue to emphasize ratio-study analysis, stratified evaluation across relevant submarkets or property classes, and careful sales verification as the foundation for evaluating mass-appraisal performance \cite{IAAO2013Ratio,IAAO2020SalesVerification,IAAO2025Mass}. Within this framework, the coefficient of dispersion (COD) remains a standard measure of uniformity, while the price-related differential (PRD) and price-related bias (PRB) remain central diagnostics for vertical equity \cite{IAAO2013Ratio,IAAO2025Mass}.

At the same time, the recent assessor literature has become more cautious about relying on any single measure of vertical equity. \textcite{IAAO2023VerticalReview} compare the behavior of several vertical-equity diagnostics and recommend reporting a \emph{suite} of tests rather than privileging one statistic. \textcite{McMillenSingh2023} further show that some regression-based procedures for detecting regressivity can be severely biased when assessments themselves are generated from regression-based valuation models, and they argue for complementing classical statistics such as the PRD with broader distribution-based diagnostics. This concern is already influencing standards discussions: the September 2025 IAAO exposure draft introduces the Vertical Equity Indicator (VEI) as an emerging principal test while still emphasizing that no single diagnostic can detect all forms of vertical inequity \cite{IAAO2025ExposureRatio}.

Beyond diagnostics, recent policy-oriented work shows that regressivity is not an isolated local anomaly but a widespread and substantively important feature of property taxation \cite{Berry2021,AvenancioLeonHoward2022}. Using millions of residential sales, \textcite{Berry2021} finds that assessments are typically regressive, that this pattern is widespread across U.S.\ counties, and that it is more consistent with limitations in assessors' data and models than with simple measurement error or explicit policy choices alone. \textcite{AvenancioLeonHoward2022} broaden the discussion by documenting a nationwide assessment gap and showing that assessment inequities can translate into systematic differences in tax burdens for the same bundle of local public services. Assessor-oriented studies also emphasize that institutional design and empirical design matter: \textcite{ScottDaigle2022} document a program that reduced both vertical and horizontal inequity, while \textcite{Rakow2023} and \textcite{Rakow2024} show that measured local regressivity and its effective incidence can depend on assessment timing, exemptions, and policy design.

A smaller but especially relevant strand of the assessor literature moves from equity measurement to equity correction within the valuation pipeline itself. \textcite{Quintos2014} propose a two-step spatial-lag/spatial-quantile procedure that improves both horizontal and vertical equity in a heterogeneous housing market. \textcite{BidansetEtAl2019} show that geographically weighted regression can reveal local vertical and horizontal inequities that are obscured by jurisdiction-wide conditional-mean models, and \textcite{BidansetEtAl2025} extend this line by using geographically weighted regression and clustering to construct more coherent market areas for mass valuation. Related work has also shifted attention from model fitting to model governance: \textcite{HermansEtAl2022} incorporate property-level ratio studies into AVM validation, and \textcite{HermansEtAl2025} use Monte Carlo sampling to assess the probability that an AVM will fail ratio-study standards under alternative sampling configurations. At the research frontier, \textcite{CandoganHanLu2023} explicitly study the accuracy--fairness trade-off in regressive property taxation using segmented, fairness-aware valuation models. Taken together, this literature shows that equity in mass appraisal does not arise automatically from better prediction alone; rather, it increasingly requires explicit measurement, explicit monitoring, and, in some settings, explicit correction within the modeling and calibration process. For the present paper, the key implication is that vertical inequity is already understood by assessors as a systematic pattern in the relationship between assessment outcomes and value level, which makes it a natural target for methods that explicitly control such relationships during training.

\subsection{Fairness in ML for Regression}

Most of the fairness-in-ML literature was originally developed for classification problems, but many real applications involve continuous prediction rather than binary decisions \cite{CatonHaas2024,BarocasHardtNarayanan2023}. In those settings, fairness cannot usually be reduced to a single thresholded decision without discarding information about score magnitudes, residual structure, and calibration. For that reason, a substantial literature has developed around \emph{fair regression}: methods that treat fairness as an explicit objective when the prediction target is continuous \cite{AgarwalDudikWu2019,FitzsimonsEtAl2019,ChzhenEtAl2020Wasserstein,ChzhenEtAl2020Recalibration,WeiEtAl2023}.

A recurring conclusion of this literature is that no single fairness notion is universally preferred for regression. Depending on the application, one may want parity in the marginal distribution of predictions, parity of losses or residuals across groups, equality of conditional means, or some broader structural notion of non-discrimination \cite{AgarwalDudikWu2019,CatonHaas2024,BarocasHardtNarayanan2023}. Methodologically, the state of the art has moved from ad hoc heuristics toward constrained or regularized empirical risk minimization. \textcite{AgarwalDudikWu2019} formulate fair regression under statistical-parity and bounded-group-loss constraints and provide reduction-based algorithms with guarantees for broad classes of losses. \textcite{FitzsimonsEtAl2019} introduce a general framework based on group fairness in expectation that can be integrated into a variety of regression learners. \textcite{WeiEtAl2023} develop an RKHS formulation for mean-parity fair regression with multiple sensitive attributes. A complementary line studies distributional fairness through optimal transport and post-processing: \textcite{ChzhenEtAl2020Wasserstein} characterize demographic-parity regression via Wasserstein barycenters, while \textcite{ChzhenEtAl2020Recalibration} derive a plug-in-and-recalibration approach with statistical guarantees. Regression-specific optimization work also includes explicit dependence constraints: \textcite{KomiyamaEtAl2018} measure unfairness through correlation or coefficient-of-determination constraints, and \textcite{ScutariPaneroProissl2022} show that fairness can be controlled through a simple ridge-type penalization strategy. Taken together, this literature establishes a useful perspective for the present paper: in continuous prediction problems, fairness can be built directly into the training or calibration process rather than being treated only as an ex-post diagnostic.

\subsection{Dependence-Based Fairness Regularization and Fair Regression with Continuous Attributes}

A particularly relevant idea from fair regression is that unfairness can often be represented as an undesirable statistical relationship between model outputs and some variable of concern; mitigation then proceeds by explicitly penalizing that relationship during training. For the purposes of the present paper, the most relevant strand is therefore the one that interprets fairness as a problem of controlling statistical dependence between predictions and sensitive variables---or, more generally, between predictions and variables along which undesirable dependence is to be controlled. This perspective is especially important when the relevant source of unfairness is not naturally binary or categorical.

Several lines of work develop this idea in different ways. \textcite{MaryEtAl2019} extend fairness-aware learning to continuous sensitive attributes and treatments using the R\'enyi maximal correlation coefficient as a differentiable dependence measure. \textcite{KomiyamaEtAl2018} formulate regression fairness directly through correlation and coefficient-of-determination constraints. \textcite{lee2022maximal} further develop maximal-correlation-based regularization as a general framework for imposing fairness criteria through dependence control. In parallel, \textcite{PerezSuayEtAl2017} and \textcite{LiEtAl2022Fairness} use kernel-based dependence measures, especially HSIC and related cross-covariance operators, to build tractable dependence-regularized learners, and \textcite{BuylDefranceDeBie2024FAIRRET} provide a recent modular framework for differentiable fairness regularization terms.

This literature is especially relevant here for two reasons. First, dependence-based regularizers are often attractive in practice because they can be incorporated into standard optimization pipelines with relatively limited changes to the underlying learner \cite{KomiyamaEtAl2018,MaryEtAl2019,LiEtAl2022Fairness,lee2022maximal,BuylDefranceDeBie2024FAIRRET}. Second, they naturally accommodate settings in which the relevant source of inequity is encoded through a continuous or structurally important variable rather than a classical protected demographic attribute. In mass appraisal, the relevant dependence is not primarily between predictions and a demographic group indicator, but between assessment behavior and value level. This makes dependence-based fair-regression methods methodologically closer to the present problem than most of the better-known group-fairness methods designed for discrete protected classes. In that sense, fair-regression theory offers assessors not just a new vocabulary, but also a practical toolkit for turning an observed inequity pattern into an explicit modeling target.

\subsection{Fairness--Accuracy Tradeoffs and Multi-Objective Model Selection}

A separate but closely related literature treats fairness mitigation as a genuinely multi-objective problem rather than as a single constrained-learning problem with one fixed penalty parameter. At a high level, this view emphasizes that different models, penalties, and post-processing choices trace out different combinations of predictive performance and fairness, so model selection should often be understood as choosing among tradeoffs rather than identifying a single universally best model \cite{CatonHaas2024}. \textcite{martinez2020minimax} formulate fairness explicitly as a multi-objective optimization problem and characterize Pareto-efficient solutions that avoid unnecessary harm across groups. \textcite{little2022fairness} study the empirical fairness--accuracy Pareto frontier directly, proposing tools to identify, quantify, and optimize model tradeoffs through frontier-based summaries rather than one-number comparisons.

This perspective is especially natural for the present application. In assessor workflows, model calibration already requires comparing many candidate specifications across multiple out-of-sample criteria, such as accuracy, dispersion, and jurisdiction-specific diagnostics. Fairness-aware mass appraisal does not create a wholly new decision problem; rather, it extends an existing calibration problem by adding equity as an additional objective. The relevant methodological question is therefore not only how to construct a regressivity-aware learner, but also how to calibrate and select candidate models when accuracy and equity move in different directions. The Pareto and frontier literature provides the natural OR language for this problem and clarifies why model selection should itself be treated as part of the fairness-aware workflow.

\subsection{Positioning of the Present Work}

The present paper sits at the intersection of these two bodies of work. On the one hand, the mass-appraisal literature has developed rich standards, diagnostics, and application-specific corrections for horizontal and vertical equity, but it still treats regressivity primarily as a ratio-study outcome or as a problem to be corrected through specific spatial, segmented, or post hoc procedures \cite{Quintos2014,IAAO2023VerticalReview,BidansetEtAl2019,CandoganHanLu2023}. On the other hand, the fair-regression literature has developed general constrained and regularized approaches for continuous prediction, including dependence-based penalties and fairness--accuracy frontier methods, but typically in settings defined by demographic protected attributes rather than assessor-style notions of vertical inequity \cite{AgarwalDudikWu2019,LiEtAl2022Fairness,lee2022maximal,martinez2020minimax,little2022fairness}.

Our contribution is to connect these two literatures in an assessor-compatible way. We start from an assessor-defined problem---regressivity as a systematic pattern in the relationship between assessment ratios and value level---and recast it in a form that fair-regression tools can act on directly during training. This shifts the role of vertical-equity analysis from purely ex-post diagnosis toward in-processing control. Concretely, we use this interpretation to build a regularization approach for standard regression models used in mass appraisal and embed the resulting corrections in a multi-objective calibration framework based on accuracy--equity tradeoffs. In this sense, the paper does not introduce equity concerns into mass appraisal, nor does it introduce fairness regularization into machine learning; rather, it develops a practical bridge between assessor-defined vertical equity and fairness-aware regression methods.

% \hrule
% \newpage
% \section{Related Work (OLD)}

% \subsection{ML in Mass Appraisal and Assessor Workflows}

% Recent reviews show that mass appraisal and real-estate valuation research has moved steadily toward automated valuation models (AVMs) and machine learning (ML) methods, with multiple regression, tree ensembles, and neural-network-based models among the most frequently studied approaches \cite{WangLi2019,Rootetal2023,Zilli2024,purdueAssessingTechniques}. Across this literature, model evaluation is still dominated by predictive criteria---typically accuracy and model comparison across algorithms---rather than by equity as an explicit model-design target \cite{WangLi2019,Rootetal2023,Zilli2024}. In particular, recent reviews identify Random Forest, Gradient Boosting, and related ensemble methods as some of the most commonly used and best-performing model classes in modern mass appraisal studies \cite{Rootetal2023,Sevgen2024100111,Zilli2024}.

% From the assessor perspective, however, adoption depends on more than predictive fit alone. The assessor-oriented AI literature emphasizes explainability, defensibility, administrative feasibility, and integration into office workflows as central conditions for real-world use \cite{IAAOAITaskForce2022}. Consistent with this, \textcite{BidansetRakow2022} report that AVM adoption in assessment offices remains partial and uneven, with perceived barriers including communication challenges, technical capacity, and budget constraints. As a result, the relevant related work for the present paper is not only the literature on which models predict best, but also the literature on how advanced models can be incorporated into assessor practice in a way that remains operationally useful and publicly defensible \cite{IAAOAITaskForce2022,BidansetRakow2022}.


% \subsection{Equity in Mass Appraisal: Horizontal Equity, Vertical Equity, and Regressivity}

% Equity has long been treated as a core objective of mass appraisal rather than as a secondary byproduct of predictive performance \cite{IAAO2013Ratio,IAAO2025Mass,Quintos2014}. In assessor practice, this objective is commonly decomposed into \emph{horizontal equity}, which concerns the uniform treatment of similarly situated properties, and \emph{vertical equity}, which concerns the absence of systematic price-related bias across the value distribution \cite{10578d49-78f8-3ef9-8277-9b6add2bd071,cornia2006property,McMillenSingh2023}. The basic diagnostic framework remains the sales-ratio study: current IAAO standards continue to emphasize ratio-study analysis, stratified evaluation across relevant submarkets or property classes, and careful sales verification as the foundation for evaluating mass-appraisal performance \cite{IAAO2013Ratio,IAAO2020SalesVerification,IAAO2025Mass}. Within this framework, the coefficient of dispersion (COD) remains a standard measure of uniformity, while the price-related differential (PRD) and price-related bias (PRB) remain central diagnostics for vertical equity \cite{IAAO2013Ratio,IAAO2025Mass}.

% At the same time, the recent assessor literature has become more cautious about relying on any single measure of vertical equity. \textcite{IAAO2023VerticalReview} compare the behavior of several vertical-equity diagnostics and recommend reporting a \emph{suite} of tests rather than privileging one statistic. \textcite{McMillenSingh2023} further show that some regression-based procedures for detecting regressivity can be severely biased when assessments themselves are generated from regression-based valuation models, and they argue for complementing classical statistics such as the PRD with broader distribution-based diagnostics. This concern is already influencing standards discussions: the September 2025 IAAO exposure draft introduces the Vertical Equity Indicator (VEI) as an emerging principal test while still emphasizing that no single diagnostic can detect all forms of vertical inequity \cite{IAAO2025ExposureRatio}.

% Beyond diagnostics, recent policy-oriented work shows that regressivity is not an isolated local anomaly but a widespread and substantively important feature of property taxation \cite{Berry2021,AvenancioLeonHoward2022}. Using millions of residential sales, \textcite{Berry2021} finds that assessments are typically regressive, that this pattern is widespread across U.S.\ counties, and that it is more consistent with limitations in assessors' data and models than with simple measurement error or explicit policy choices alone. \textcite{AvenancioLeonHoward2022} broaden the discussion by documenting a nationwide assessment gap and showing that assessment inequities can translate into systematic differences in tax burdens for the same bundle of local public services. At the same time, assessor-oriented studies emphasize that institutional design and empirical design both matter: \textcite{ScottDaigle2022} document a program that reduced both vertical and horizontal inequity, while \textcite{Rakow2023} and \textcite{Rakow2024} show that measured local regressivity and its effective incidence can depend on assessment timing, exemptions, and policy design.

% A smaller but especially relevant strand of the assessor literature moves from equity measurement to equity correction within the valuation pipeline itself. \textcite{Quintos2014} propose a two-step spatial-lag/spatial-quantile procedure that improves both horizontal and vertical equity in a heterogeneous housing market. \textcite{BidansetEtAl2019} show that geographically weighted regression can reveal local vertical and horizontal inequities that are obscured by jurisdiction-wide conditional-mean models, and \textcite{BidansetEtAl2025} extend this line by using geographically weighted regression and clustering to construct more coherent market areas for mass valuation. Related work has also shifted attention from model fitting to model governance: \textcite{HermansEtAl2022} incorporate property-level ratio studies into AVM validation, and \textcite{HermansEtAl2025} use Monte Carlo sampling to assess the probability that an AVM will fail ratio-study standards under alternative sampling configurations. At the research frontier, \textcite{CandoganHanLu2023} explicitly study the accuracy--fairness trade-off in regressive property taxation using segmented, fairness-aware valuation models. Taken together, this literature shows that equity in mass appraisal does not arise automatically from better prediction alone; rather, it increasingly requires explicit measurement, explicit monitoring, and, in some settings, explicit correction within the modeling and calibration process.

% \subsection{Fairness in ML for Regression}

% Fairness in machine learning has matured into a broad literature spanning group, individual, and causal notions of non-discrimination, but many of its canonical definitions and mitigation methods were originally developed for classification rather than continuous prediction \cite{CatonHaas2024,BarocasHardtNarayanan2023,DworkEtAl2012,KusnerEtAl2017}. Regression settings raise additional issues because thresholding continuous outputs generally discards information about score magnitudes, residual structure, and calibration. As a result, the fair-regression literature has developed continuous-output analogues of fairness notions rather than simply importing classification metrics unchanged \cite{AgarwalDudikWu2019,FitzsimonsEtAl2019,ChzhenEtAl2020Wasserstein,ChzhenEtAl2020Recalibration,WeiEtAl2023}. A recurring conclusion of this literature is that no single fairness notion is universally preferred for regression: the appropriate criterion depends on whether the application seeks parity in the marginal distribution of predictions, parity of losses or residuals across groups, equality of conditional means, or a more structural notion of non-discrimination \cite{AgarwalDudikWu2019,CatonHaas2024,BarocasHardtNarayanan2023}.

% Methodologically, an important development in fair regression has been the move from model-specific heuristics to general constrained or regularized empirical risk minimization. \textcite{AgarwalDudikWu2019} formulate fair regression under statistical-parity and bounded-group-loss constraints and provide reduction-based algorithms with guarantees for broad classes of loss functions. \textcite{FitzsimonsEtAl2019} introduce a general framework based on group fairness in expectation that can be integrated into a variety of regression learners. \textcite{WeiEtAl2023} develop an RKHS formulation for mean-parity fair regression with multiple sensitive attributes. A complementary line studies distributional fairness through optimal transport and post-processing: \textcite{ChzhenEtAl2020Wasserstein} characterize demographic-parity regression via Wasserstein barycenters, while \textcite{ChzhenEtAl2020Recalibration} derive a plug-in-and-recalibration approach with statistical guarantees. Regression-specific optimization work also includes explicit dependence constraints: \textcite{KomiyamaEtAl2018} measure unfairness through correlation or coefficient-of-determination constraints, and \textcite{ScutariPaneroProissl2022} show that fairness can be controlled through a simple ridge-type penalization strategy. Together, these papers establish a clear view of fair regression as a setting in which fairness is itself an explicit modeling target rather than merely an ex-post diagnostic.

% \subsection{Dependence-Based Fairness Regularization and Fair Regression with Continuous Attributes}

% For the purposes of the present paper, the most relevant fair-regression strand is the one that interprets fairness as a problem of controlling statistical dependence between predictions and sensitive variables, or between suitable transformations of them. This perspective is particularly important when the relevant source of unfairness is not naturally binary or categorical. \textcite{MaryEtAl2019} extend fairness-aware learning to continuous sensitive attributes and treatments using the R\'enyi maximal correlation coefficient as a differentiable dependence measure. \textcite{KomiyamaEtAl2018} formulate regression fairness directly through correlation and coefficient-of-determination constraints. \textcite{lee2022maximal} further develop maximal-correlation-based regularization as a general framework for imposing fairness criteria through dependence control. In parallel, \textcite{PerezSuayEtAl2017} and \textcite{LiEtAl2022Fairness} use kernel-based dependence measures, especially HSIC and related cross-covariance operators, to build tractable dependence-regularized learners, and \textcite{BuylDefranceDeBie2024FAIRRET} provide a recent modular framework for differentiable fairness regularization terms.

% This literature is especially relevant here for two reasons. First, dependence-based regularizers are computationally convenient surrogates for broader independence objectives and can often be incorporated into standard optimization pipelines with relatively limited changes to the underlying learner \cite{KomiyamaEtAl2018,MaryEtAl2019,LiEtAl2022Fairness,lee2022maximal,BuylDefranceDeBie2024FAIRRET}. Second, they naturally accommodate settings in which the relevant source of inequity is encoded through a continuous or structurally important variable rather than a classical protected demographic attribute. In mass appraisal, the relevant dependence is not primarily between predictions and a demographic group indicator, but between assessment behavior and value level. This makes dependence-based fair-regression methods methodologically closer to the present problem than most of the better-known group-fairness methods designed for discrete protected classes.

% \subsection{Fairness--Accuracy Tradeoffs and Multi-Objective Model Selection}

% A separate but closely related literature treats fairness mitigation as a genuinely multi-objective problem rather than as a single constrained-learning problem with one fixed penalty parameter. At a high level, this view emphasizes that different models, penalties, and post-processing choices trace out different combinations of predictive performance and fairness, so model selection should often be understood as choosing among tradeoffs rather than identifying a single universally best model \cite{CatonHaas2024}. \textcite{martinez2020minimax} formulate fairness explicitly as a multi-objective optimization problem and characterize Pareto-efficient solutions that avoid unnecessary harm across groups. \textcite{little2022fairness} study the empirical fairness--accuracy Pareto frontier directly, proposing tools to identify, quantify, and optimize model tradeoffs through frontier-based summaries rather than one-number comparisons.

% This perspective is especially natural for the present application. In assessor workflows, model calibration already requires comparing many candidate specifications across multiple out-of-sample criteria, and fairness-aware mass appraisal adds at least one more objective to that process. The relevant methodological question is therefore not only how to construct a regressivity-aware learner, but also how to calibrate and select candidate models when accuracy and equity move in different directions. The Pareto and frontier literature provides the natural OR language for this problem.

% \subsection{Positioning of the Present Work}

% The present paper sits at the intersection of these two bodies of work. On the one hand, the mass-appraisal literature has developed rich standards, diagnostics, and application-specific corrections for horizontal and vertical equity, but it still treats regressivity primarily as a ratio-study outcome or as a problem to be corrected through specific spatial, segmented, or post hoc procedures \cite{Quintos2014,IAAO2023VerticalReview,BidansetEtAl2019,CandoganHanLu2023}. On the other hand, the fair-regression literature has developed general constrained and regularized approaches for continuous prediction, including dependence-based penalties and fairness--accuracy frontier methods, but typically in settings defined by demographic protected attributes rather than assessor-style notions of vertical inequity \cite{AgarwalDudikWu2019,LiEtAl2022Fairness,lee2022maximal,martinez2020minimax,little2022fairness}.

% Our contribution is to connect these two literatures in an assessor-compatible way. We treat regressivity as a dependence problem between assessment ratios and value level, use this interpretation to build an in-processing regularization approach for standard regression models used in mass appraisal, and embed the resulting corrections in a multi-objective calibration framework based on accuracy--equity tradeoffs. In this sense, the paper does not introduce equity concerns into mass appraisal, nor does it introduce fairness regularization into machine learning; rather, it develops a practical bridge between assessor-defined vertical equity and fairness-aware regression methods.

% % OLD RELATED WORK SECTION
% % \hrule
% % \newpage
% % \section{Related Work: Maybe just use the introduction, since there not much about the specific topic (regressivity in ML-based MA)}

% % \subsection{Equity in Mass Appraisal}

% % Equity is a core objective of mass appraisal, not a secondary byproduct of predictive accuracy, in both the assessor standards literature and the broader applied literature on assessment performance \cite{IAAO2013Ratio,IAAO2025Mass,Quintos2014}. In assessor practice, this objective is usually decomposed into \emph{horizontal equity}, which concerns the uniform treatment of similarly situated properties or property groups, and \emph{vertical equity}, which concerns the absence of systematic price-related bias across the value distribution; in practice, the latter is most often discussed through the lens of \emph{regressivity}, whereby lower-valued properties are assessed at higher fractions of market value than higher-valued properties \cite{IAAO2013Ratio,IAAO2025Mass,McMillenSingh2023}. Consistent with this framing, the current \textcite{IAAO2025Mass} standard emphasizes that assessment uniformity must be examined not only overall, but also across neighborhoods, construction classes, age groups, and size groups, and it continues to treat the sales-ratio study as the basic diagnostic framework for evaluating mass appraisal performance \cite{IAAO2025Mass}. The same standards literature also stresses that equity diagnostics are inseparable from data quality: \textcite{IAAO2020SalesVerification} states that only sales meeting the definition of market value and properly adjusted for non-realty components should be used in valuation and ratio-study work, and that appraisal accuracy depends on proper verification and adjustment of sales data \cite{IAAO2020SalesVerification}. Within this framework, the classical toolkit evaluates within-group uniformity primarily through the coefficient of dispersion (COD) and price-related equity primarily through the price-related differential (PRD) and price-related bias (PRB), with the IAAO standards continuing to interpret PRDs above 1.03 as tending to indicate regressivity and PRDs below 0.98 as tending to indicate progressivity \cite{IAAO2013Ratio,IAAO2025Mass}. Recent literature, however, has become much more cautious about relying on any single vertical-equity statistic: \textcite{IAAO2023VerticalReview} compare the behavior of competing vertical-equity measures under simulation and explicitly recommend reporting a \emph{suite} of tests rather than privileging one measure, while \textcite{McMillenSingh2023} show that regression-based procedures for detecting regressivity are severely biased when assessments themselves are generated from regression-based valuation models and argue that the PRD should be complemented by distribution-based procedures such as Gini-, Suits-, and kernel-density-based diagnostics \todo{do we connect this part with the higher density measures like dCorr, maximal corr., etc.?} \cite{IAAO2023VerticalReview,McMillenSingh2023}. This reassessment of measurement practice is already influencing the standards conversation: the September 2025 IAAO exposure draft introduces the Vertical Equity Indicator (VEI) as an emerging principal test, while still acknowledging that no single test can detect all vertical inequities and that multiple diagnostics remain necessary \cite{IAAO2025ExposureRatio}.

% % % \subsubsection{Equity in Assessment: Policy View}

% % \textbf{The policy point of view.} On the substantive side, the policy literature now documents that regressivity is not an isolated local anomaly, but a widespread empirical feature of property taxation with meaningful distributive consequences \cite{Berry2021,AvenancioLeonHoward2022}. Using millions of residential sales, \textcite{Berry2021} finds that assessments are typically regressive, that low-priced properties are assessed at higher fractions of sale price than high-priced properties, that this pattern is present in at least two-thirds of U.S.\ counties, and that it is more consistent with limitations in assessors' data and models than with simple measurement error or explicit policy choices alone \cite{Berry2021}. \textcite{AvenancioLeonHoward2022} broaden the discussion by documenting a nationwide ``assessment gap,'' showing that, holding jurisdictions and tax rates fixed, Black and Hispanic residents face roughly 10--13\% higher property-tax burdens for the same bundle of public services, and tracing this disparity in part to assessments being less sensitive to neighborhood attributes than market prices and to racial differences in appeals behavior and appeals outcomes within neighborhoods \cite{AvenancioLeonHoward2022}. Assessor-oriented responses and follow-up studies accept the importance of these findings while emphasizing that institutional design and empirical design both matter: \textcite{ScottDaigle2022} find that New Jersey's Assessment Demonstration Program reduced both vertical and horizontal inequity in Monmouth County, whereas \textcite{Rakow2023} argues that empirical conclusions about local regressivity can be sensitive to the alignment of sales years and assessment years and to whether homestead exemptions are incorporated, and \textcite{Rakow2024} further argues that flat-dollar homestead exemptions can partially offset regressive tax incidence when assessment systems remain imperfect \cite{ScottDaigle2022,Rakow2023,Rakow2024}. 

% % \textbf{Equity Correction in ML-based Mass Appraisal.}
% % At the same time, the recent machine-learning and AVM literature shows that improvements in predictive fit do not automatically resolve these equity problems \cite{WangLi2019,Zilli2024,IAAOAITaskForce2022}. Broad reviews of mass appraisal models report that recent work is dominated by accuracy-oriented questions and that Random Forest and Gradient Boosting are among the most frequently used and strongest-performing models, but they also note substantial heterogeneity in performance metrics and relatively limited attention to equity as a primary model-design target \cite{WangLi2019,Zilli2024}. This gap is especially salient from the assessor perspective: \textcite{IAAOAITaskForce2022} review AI methods for assessment offices while emphasizing organizational impacts, model transparency, and administrative feasibility, and \textcite{BidansetRakow2022} report that only 16\% of surveyed offices were using regression-based AVMs even though respondents associated them with gains in accuracy, uniformity, equity, and efficiency, with non-adoption driven largely by perceived adequacy of current methods, difficulty explaining results to taxpayers and staff, and budget and technical constraints \cite{IAAOAITaskForce2022,BidansetRakow2022}. The part of the literature that most directly addresses equity therefore modifies the valuation pipeline itself rather than assuming equity will emerge from better prediction alone. In the assessor literature, \textcite{Quintos2014} is a particularly important example: using single-family homes in Brooklyn, she proposes a two-step spatial-lag/spatial-quantile framework in which the spatial step improves COD and addresses heterogeneity, while the quantile step reduces the correlation between assessment ratios and price from $-44.05\%$ to $-8.22\%$, yields a PRD of 1.0208 that satisfies IAAO guidance, and concludes that the spatial-quantile approach can correct both horizontal and vertical equity in a highly heterogeneous market \cite{Quintos2014}. Subsequent spatial work by \textcite{BidansetEtAl2019} shows that geographically weighted regression (GWR) can reveal vertical and horizontal inequities at the intra-neighborhood level that are obscured by jurisdiction-wide conditional-mean models, and \textcite{BidansetEtAl2025} extend this line by proposing automated market-area delineation using GWR and hierarchical clustering, reporting predictions that satisfy IAAO accuracy and uniformity thresholds both citywide and within the proposed market areas \cite{BidansetEtAl2019,BidansetEtAl2025}. Related assessor work has also moved from model fitting to model governance by incorporating property-level ratio studies into AVM validation workflows and by using Monte Carlo sampling to estimate the probability that an AVM will fail to comply with ratio-study standards under alternative sample configurations \cite{HermansEtAl2022,HermansEtAl2025}. At the research frontier, \textcite{CandoganHanLu2023} develop a segmented, fairness-aware property valuation model that explicitly studies the accuracy--fairness trade-off in regressive property taxation, illustrating the broader shift toward models that target equity directly rather than treating it as an incidental validation outcome \cite{CandoganHanLu2023}. Taken together, the current literature suggests that equity in mass appraisal has become a genuinely multi-layered problem: horizontal equity is closely tied to local heterogeneity, market-area delineation, and within-group dispersion, whereas vertical equity is closely tied to systematic price-related bias and therefore requires explicit measurement, explicit monitoring, and, increasingly, explicit correction within the modeling and calibration pipeline \cite{IAAO2025Mass,IAAO2023VerticalReview,Berry2021,Quintos2014,BidansetEtAl2025,CandoganHanLu2023}.

% % \hrule


% % % \subsection{Equity in Mass Appraisal}


% % % Equity is a foundational objective of mass appraisal, not merely a secondary diagnostic of predictive performance. In the assessor literature, this objective is usually decomposed into \textit{horizontal equity}, which concerns the uniform treatment of similarly situated properties or property groups, and \textit{vertical equity}, which concerns the absence of systematic price-related bias across the value distribution; in practice, the latter is most often discussed through the lens of \textit{regressivity}, whereby lower-valued properties are assessed at higher fractions of market value than higher-valued properties. Current IAAO guidance still treats \textit{sales-ratio analysis} as the indispensable “bottom-line” evaluation of appraisal performance regardless of how values are generated, and it emphasizes that assessment uniformity should be examined not only overall but also across neighborhoods, construction classes, age groups, size groups, and other relevant strata. This literature also makes clear that equity measurement is inseparable from data quality: sales used for model calibration and ratio studies must be verified, adjusted where needed, and screened for representativeness, because inaccurate or non-arm’s-length sales contaminate both the valuation model and the subsequent diagnosis of horizontal and vertical inequity. 

% % % Within this framework, the most developed recent debate concerns the measurement of vertical equity. The traditional assessor toolkit relies on the Coefficient of dispersion (COD) as a measure of dispersion or uniformity and on the Price-related differential (PRD) and Price-related Bias (PRB) as summary measures of price-related bias. Under the still-operative 2013 Standard on Ratio Studies, PRDs above 1.03 are interpreted as tending to indicate regressivity and PRDs below 0.98 as tending to indicate progressivity. However, the recent literature has become much more cautious about relying on any single vertical-equity statistic. The IAAO Statistical Tools and Measures Task Force (2023) explicitly recommends reporting a ``suite of tests'' rather than privileging one measure, while McMillen and Singh (2023) show that many commonly used regression-based procedures for detecting regressivity are severely biased, especially when assessments themselves are produced from regression-based valuation models. Their analysis therefore both preserves a role for familiar summary measures such as the PRD and motivates complementary distributional or nonparametric diagnostics. This concern is now influential enough that IAAO’s 2025 exposure draft introduces the Vertical Equity Indicator (VEI) as an emerging main test, yet the same draft also reiterates that no single test should be expected to detect all vertical-equity problems and that the 2013 standard remains the official guidance until a final revision is adopted. Taken together, the current assessor literature no longer treats vertical equity as a one-metric problem; instead, it increasingly views it as a multi-diagnostic question that requires robustness to outliers, nonlinearities, model form, and sample composition. 

% % % On the substantive side, recent policy-oriented work has renewed attention to regressivity as a widespread and economically meaningful feature of property taxation. Berry (2021), using millions of residential transactions from more than 2,600 counties, finds that assessments are typically regressive, that such regressivity is present in at least two-thirds of counties, and that it disproportionately burdens owners of lower-valued homes; importantly, Berry argues that the pattern cannot be explained away by sales-price measurement error or by explicit policy devices such as assessment limits alone, but instead appears to arise from limitations in assessors’ data and models. Avenancio-León and Howard (2022) broaden the discussion by documenting a nationwide “assessment gap,” under which Black and Hispanic residents face roughly 10–13\% higher property-tax burdens for the same bundle of local public services; they trace this disparity in part to assessments being less sensitive to neighborhood attributes than market prices and to differences in appeals behavior and appeals outcomes within neighborhoods. At the same time, assessor-facing policy studies emphasize that institutional design can mitigate, though not eliminate, these inequities: Scott and Daigle (2022) show that New Jersey’s Assessment Demonstration Program reduced both vertical and horizontal inequity in Monmouth County, while Rakow (2023, 2024) argues that properly structured homestead exemptions can partially offset regressive tax incidence even when assessment systems remain imperfect. Thus, the policy literature has two consistent messages: first, inequity in assessments is real and consequential; second, it reflects an interaction between valuation models, reassessment practices, appeals, neighborhood effects, and relief policy rather than a single mechanical source. 

% % % Against this backdrop, the recent ML and Automated Valuation Models (AVM) literature reveals a clear tension between \textit{accuracy gains} and \textit{equity control}. Broad reviews of mass appraisal show that recent research is dominated by predictive questions, with Random Forest and Gradient Boosting appearing especially frequently and often outperforming traditional appraisal approaches on accuracy metrics; however, these same reviews also note substantial heterogeneity in performance measures and relatively limited attention to equity as a primary endpoint. The assessor literature on AI adoption is even more explicit that accuracy alone is insufficient: the IAAO AI Task Force argues that AI adoption in assessment offices must be: (i) explainable, (ii) defensible, and (iii) trusted by both staff and taxpayers, while Bidanset and Rakow’s survey reports that regression-based AVMs are still used by only a minority of offices and that barriers include explainability, budget, and technical expertise, even though offices using AVMs report gains in accuracy, uniformity, and equity. As a result, the most relevant methodological contributions for the present literature review come from a smaller strand of work that explicitly modifies the valuation pipeline to address inequity. Quintos (2014) is especially important because it treats horizontal and vertical equity jointly: in her two-step spatial-lag/spatial-quantile procedure, the spatial step improves COD and addresses heterogeneity and heteroskedasticity, while the quantile step materially flattens the ratio-price relationship, reducing the absolute correlation between assessment ratio and price from 44.05\% to 8.22\% while keeping COD within IAAO bounds. Related spatial work by Bidanset and coauthors argues that geographically weighted regression can uncover intra-neighborhood patterns of both horizontal and vertical inequity and, more recently, can also be used with clustering methods to automate coherent market-area delineation for tax assessment. More recent assessor literature has also extended the discussion from model fitting to model governance through composite modeling and Monte Carlo-based quality assurance, while the broader fairness-in-ML frontier is represented by Candogan, Han, and Lu’s 2023 preprint, which explicitly optimizes the accuracy-fairness trade-off in regressive property taxation through segmented models. The overarching implication of this literature is that equity does not arise automatically from better fit: horizontal equity generally requires improved local or spatial representation of heterogeneous markets, whereas vertical equity requires explicit monitoring and, increasingly, explicit correction of price-related bias. ([oaj.fupress.net][1])



% % % % [1]: https://oaj.fupress.net/index.php/ceset/article/view/15792?utm_source=chatgpt.com "Machine learning models in mass appraisal for property tax ..."





% % \subsection{Machine Learning in Mass Appraisal}

% % There is a lot of this. Select the best ones probably.

% % \textbf{Case Study.} 
% % There exists a well-documented example of the use of advanced ML models in the mass appraisal process: the Cook County Assessor's Office (CCAO) in Chicago \cite{ccao_how_properties_are_valued}. The CCAO manages an extensive open-access repository with the datasets, code, and guides describing the mass appraisal process used to predict property values in Cook County \cite{ccao_data_2026}. Linear models were used by the CCAO until 2018, and then the office transitioned to gradient boosting methods in 2019, specifically LightGBM \cite{lightgbm}. The latter methods achieve high predictive accuracy, but they also exhibit stronger regressive behavior \cite{CandoganHanLu2023}. We aim to reduce these regressive patterns while maintaining a similar level of accuracy. 



% % \subsection{Fairness in Machine Learning}

% % \begin{itemize}
% % \item Specially correlation and penalized corrections or mitigations (usually, with sensitive attributes).
% %     \item 
% % Regressivity correction for Mass Appraisal. The super-close paper for achieving less regressivity: Haihao's.
% % \end{itemize}

% % \subsection{Fairness in Machine Learning for Regression}

% % Fairness in machine learning has matured into a broad literature spanning group, individual, and causal notions of non-discrimination, but most of its canonical definitions, benchmarks, and mitigation algorithms were originally developed for binary classification rather than continuous prediction \cite{CatonHaas2024,BarocasHardtNarayanan2023}. At the conceptual level, current surveys and reference texts organize observational group fairness around independence, separation, and sufficiency, individual fairness around the principle that similar individuals should receive similar decisions, and causal fairness around counterfactual or path-specific invariance \cite{BarocasHardtNarayanan2023,DworkEtAl2012,KusnerEtAl2017}. For regression, however, these notions cannot in general be reduced to a single thresholded decision without discarding information about score magnitudes, residual structure, or calibration, so the fair-regression literature has developed continuous-output analogues of these ideas rather than simply importing classification metrics unchanged \cite{AgarwalDudikWu2019,ChzhenEtAl2020Wasserstein,ChzhenEtAl2020Recalibration,FitzsimonsEtAl2019,WeiEtAl2023}. A recurring conclusion of this literature is that there is no universally preferred fairness metric for regression: the appropriate notion depends on whether the application seeks parity in the marginal distribution of predictions, parity of losses or residual behavior across groups, equality of conditional means, or a more structural notion of non-discrimination \cite{AgarwalDudikWu2019,BarocasHardtNarayanan2023,CatonHaas2024}. Surveys also continue to organize mitigation methods into pre-processing, in-processing, and post-processing workflows, but increasingly emphasize that this is not a merely taxonomic distinction: it determines whether fairness is treated as a data-repair problem, a constrained statistical learning problem, or a post hoc calibration problem \cite{CatonHaas2024}.

% % A major development in the state of the art for fair regression is the move from model-specific heuristics to general constrained empirical risk minimization. \textcite{AgarwalDudikWu2019} formulate fair regression under statistical parity and bounded group-loss constraints and give reduction-based algorithms with guarantees for arbitrary Lipschitz losses, including least-squares, logistic, and quantile-type losses. In parallel, \textcite{FitzsimonsEtAl2019} introduce group fairness in expectation as a generic constraint that can be incorporated into kernel methods, Gaussian processes, support vector machines, neural network regression, and tree-based regression, while \textcite{WeiEtAl2023} develop an RKHS framework for mean-parity fair regression that is compatible with multiple sensitive attributes and also connects naturally to covariance-based fairness notions. A complementary line studies distributional fairness through optimal transport: \textcite{ChzhenEtAl2020Wasserstein} characterize the optimal demographic-parity regressor through Wasserstein barycenters, and \textcite{ChzhenEtAl2020Recalibration} derive a two-stage plug-in-and-recalibration procedure with statistical guarantees; \textcite{TaturyanChzhenHebiri2024} extend this post-processing logic to the practically important setting in which sensitive attributes are unavailable at inference time. Regression-specific optimization approaches also include explicit low-order dependence constraints: \textcite{KomiyamaEtAl2018} define unfairness through the correlation between predictions and sensitive attributes and formulate fair regression with correlation or coefficient-of-determination constraints, while \textcite{ScutariPaneroProissl2022} show that a user-specified fairness level can be enforced through a ridge-penalized model-selection procedure that controls the effect of sensitive variables. Together, these papers establish a clear current view of fair regression as a constrained or regularized prediction problem in which fairness is an explicit modeling target rather than an incidental validation diagnostic \cite{AgarwalDudikWu2019,FitzsimonsEtAl2019,KomiyamaEtAl2018,ScutariPaneroProissl2022}.

% % For the direction of the present paper, the most relevant strand of the literature is the one that interprets fairness as a problem of controlling statistical dependence between predictions and sensitive variables, or between suitable functions of them. \textcite{MaryEtAl2019} generalize fairness-aware learning to continuous sensitive attributes and treatments by using the R\'enyi maximal correlation coefficient as a differentiable dependence measure; \textcite{PerezSuayEtAl2017} and \textcite{LiEtAl2022Fairness} use kernel-based dependence penalties, especially HSIC and related cross-covariance operators, to obtain tractable fair regression estimators and corresponding Gaussian-process interpretations; and \textcite{lee2022maximal} show that maximal-correlation regularizers can be used to enforce independence- and separation-type criteria for both discrete and continuous variables. More recently, \textcite{BuylDefranceDeBie2024FAIRRET} propose a modular differentiable regularization framework that makes a wide class of fairness penalties easier to incorporate into modern automatic-differentiation pipelines. This line of work is especially relevant because covariance-, correlation-, and kernel-dependence penalties can all be viewed as computationally convenient surrogates for broader independence objectives, while remaining compatible with gradient-based optimization and standard learners \cite{KomiyamaEtAl2018,MaryEtAl2019,PerezSuayEtAl2017,LiEtAl2022Fairness,lee2022maximal,BuylDefranceDeBie2024FAIRRET}. At the same time, the broader fair-ML literature cautions that fairness interventions are highly sensitive to the exact problem formulation, data regime, and evaluation protocol, which has motivated benchmark and frontier-based analyses rather than single-point comparisons \cite{HanEtAl2024FFB,DefranceBuylDeBie2024ABCFair}. This supports the increasingly common view that fairness mitigation should be treated as a genuinely multi-objective problem, in which the practitioner explores or optimizes an accuracy--fairness frontier rather than fixing one penalty value a priori \cite{martinez2020minimax,little2022fairness}. Finally, recent work on fairness without demographic data and fairness under distribution shift shows that deployment constraints and distributional instability are now part of the state of the art rather than edge cases \cite{AshurstWeller2023,KenfackEtAl2024,ShaoEtAl2024,TaturyanChzhenHebiri2024}. In that sense, a workflow based on dependence-aware regularization, robust penalty calibration or model selection, and a second-stage constrained correction is well aligned with one of the most active directions in current fair-regression research, even if the present application differs from the classical protected-group setting in that the ``sensitive'' direction is encoded through a structurally important continuous variable rather than a standard demographic attribute \cite{MaryEtAl2019,LiEtAl2022Fairness,KomiyamaEtAl2018,ChzhenEtAl2020Recalibration,martinez2020minimax,little2022fairness}.


% % \subsubsection{Regressivity in ML-based Mass Appraisal}

% % % \textbf{Why can ML enforce regressivity in assessments?} Let $y \in \Y$ be the log-prices (logarithm of the market price) vector, and let $f(X)=\Expected{y|X}$ be the ideal predictor (Bayes-optimal regressor) that minimizes the squared-error. Let $e = f(X)-y$ be the vector of residuals. Then
% % % \begin{align*}
% % % \Cov(e,y)= & \Cov(f(X)- y, y) = \Cov(\Expected{y|X}, y) - \Var(y) \\
% % % =& Var(\Expected{y|X}) - Var(y) = -\Expected{Var(y|X)}\\
% % % \Longrightarrow \Cov(e,y) \leq 0
% % % \end{align*}
% % % That is, as far as there exist variance of $y$ conditional on $X$ (which is specially true for real estate, where prices have a lot of noise), then the covariance between the $\log r$ (in terms of the Bayes-optimal predictor) and the log-prices $y$ is going to be negative, i.e., there is going to exist a negative slope between the two quantities.



% % \subsection{Multi-objective model selection}
% % \todo{Other point: Robustness for selecting it or guarantees for OOS Pareto frontier.}


% % Generally speaking, the methodology for selecting the best performing model and configuration of hyperparameters in out-of-sample (OOS) sets its based on Cross Validation (CV)---and its variations---for unstructured data. In the case of time-dependencies, the usual methodology is rolling-origin validation process or similar strategies. Usually, a practitioner would select the best performing model by choosing the model with the best out-of-sample average accuracy for some measure of prediction error, or some aggregation of multiple metrics. However, on the multi-objective optimization paradigm (e.g., accuracy, dispersion, fairness, etc.) different models could be the best average performing model for different metrics. For example, one model could have the best accuracy performance OOS, but perform poorly in fairness senses, while the opposite can also be true. This difference the model selection paradigm implies that, while in the single-objective case we have exactly one optimal performance measurement, in the multi-objective case we have set of Pareto-optimal solutions (infinite candidates), with of optimal performance; one for each possible tradeoff on the Pareto  frontier. A graphic representation of this difference is shown in Figure \ref{fig:model_selection}.
% % \begin{figure}[!h]
% %     \centering
% %     \caption{Difference between single-objective model selection, and multi-objective model selection. On the single-objective case, one selects a model that reaches the highest accuracy, while in the multi-objective case we could select any of the solutions in the Pareto-optimal front, which represent different optimal tradeoffs between prediction Error \& Regressivity.}
% %     \includegraphics[width=\linewidth]{Fairness Price/img/model_selection/model_selection_paradigms.pdf}
% %     \label{fig:model_selection}
% % \end{figure}









\section{Problem Formulation, Equity Measures, and Temporal Evaluation}
\label{sec:problem_setting}

This section fixes the basic prediction notation, defines the assessor-facing equity objects used throughout the paper, and formalizes the temporal evaluation setting in which models are trained, calibrated, and compared. The goal is to describe the predictive problem in the same language used by both modern regression workflows and ratio-study practice, so that the methodological section can build on a single, stable set of definitions.

\subsection{Mass-Appraisal Prediction Setting}
\label{subsec:prediction_setting}

Mass appraisal estimates the value of a \emph{group} of properties as of a given date using common data, standardized methods, and statistical testing \cite{IAAO2025Mass}. In the predictive setting studied here, we observe a sample of $n$ sold properties
$
\mathcal{D}_n := \{(x_i,p_i,t_i)\}_{i=1}^n,
$
where $x_i \in \mathcal{X}$ is the vector of observed characteristics for property $i \in \{1,\ldots,n\}$, $p_i \in \mathbb{R}_{++}$ is its verified sale price, and $t_i$ is its sale date. As in ratio-study practice, sale price is used as the market-value proxy on sold parcels.

Because prices are positive and strongly right-skewed, we work internally on the log-price scale. For each sold property, define
$
y_i := \log p_i.
$
Let $\Theta$ denote the parameter space of a regression model, and let
$
f_\theta:\mathcal{X}\to\mathbb{R},
$ with $ \theta \in \Theta,
$
denote a candidate predictor of log price. For a given property $i$, the model produces the log-scale prediction
$
\hat y_i(\theta) := f_\theta(x_i),
$
and the corresponding prediction on the original price scale is
$
\hat p_i(\theta) := \exp(\hat y_i(\theta)).
$
Throughout the paper, $\hat p_i(\theta)$ will be interpreted as the model-implied assessed value at the predictive layer. 
s
The \textit{baseline} learning problem is the regularized empirical risk minimization problem
\begin{equation}
\label{eq:baseline_erm}
\hat\theta
\in
\arg\min_{\theta\in\Theta}
% \left\{
\frac{1}{n}\sum_{i=1}^n \ell\!\big(f_\theta(x_i),y_i\big)
+
\Omega(\theta)
% \right\}
,
\end{equation}
where $\ell$ is a loss function on the log-price scale and $\Omega$ is an optional complexity penalty or model-specific regularizer. This formulation covers the main model classes studied later, including linear regression and boosting-based learners.

A central identity for the paper is obtained by comparing predictions and observed prices on the log scale. Define the log-scale residual
% \begin{equation}
% \label{eq:residual_def}
$
e_i(\theta) := \hat y_i(\theta)-y_i.
$
% \end{equation}
Then
\begin{equation}
\label{eq:log_ratio_identity}
e_i(\theta)
=
\log \hat p_i(\theta)-\log p_i
=
\log\!\left(\frac{\hat p_i(\theta)}{p_i}\right).
\end{equation}
Thus the residual is exactly the logarithm of the assessment ratio. This identity is the main formal link between the regression problem and the assessor-side ratio-study framework.

\subsection{Assessment Ratios and Equity Measures}
\label{subsec:equity_metrics}

For any fitted model and any evaluation sample $S \subseteq \{1,\dots,n\}$, define the \emph{assessment-to-sale ratio}
\begin{equation}
\label{eq:ratio_def}
r_i := \frac{\hat p_i}{p_i},
\qquad i\in S,
\end{equation}
where dependence on the fitted model is suppressed when clear from context. Ratio studies are conducted on the original price scale, so all equity diagnostics below are computed from $\hat p_i$ and $p_i$, not from log-price residuals directly.

Under proportional assessment, the ratios $\{r_i\}_{i\in S}$ should remain stable across value levels. Regressivity arises when lower-valued properties tend to have larger ratios than higher-valued properties, so that lower-end parcels are assessed at larger fractions of market value than higher-end parcels. Progressivity is the opposite pattern.

Let
\[
m_S := \operatorname{med}\{r_i: i\in S\}
\]
denote the median ratio in sample $S$. We summarize horizontal equity primarily through the coefficient of dispersion (COD), which the IAAO describes as the average percentage deviation from the median ratio \cite{IAAO2013Ratio,IAAO2025ExposureRatio}:
\begin{equation}
\label{eq:cod}
\operatorname{COD}(S)
:=
\frac{100}{|S|\,m_S}
\sum_{i\in S} |r_i-m_S|.
\end{equation}
Lower COD values indicate greater within-group uniformity.

We summarize vertical equity using price-related diagnostics. The price-related differential (PRD) is
\begin{equation}
\label{eq:prd}
\operatorname{PRD}(S)
:=
\frac{\frac{1}{|S|}\sum_{i\in S} r_i}
{\frac{\sum_{i\in S} \hat p_i}{\sum_{i\in S} p_i}},
\end{equation}
that is, the mean ratio divided by the value-weighted mean ratio. Under standard assessor interpretation, PRD values above $1.03$ tend to indicate regressivity, while values below $0.98$ tend to indicate progressivity \cite{IAAO2013Ratio,IAAO2025ExposureRatio}.

We also report the coefficient of price-related bias (PRB), computed according to the IAAO specification \cite{IAAO2025ExposureRatio}. Conceptually, the PRB is a regression-based measure of the percentage change in assessment ratios when value doubles or halves, using a value proxy designed to reduce measurement-error bias. Negative PRB values indicate regressivity, positive values indicate progressivity, and values near zero indicate approximate price proportionality.

The key point for the present paper is therefore simple: assessor-side equity is evaluated through ratio-based summaries on the original price scale, while the learning problem is solved on the log-price scale. The identity in \eqref{eq:log_ratio_identity} is what allows these two views to be connected.

\subsection{Regressivity as a Dependence Problem}
\label{subsec:dependence_problem}

The assessor literature usually describes regressivity as a \emph{price-related} pattern in the ratios: as property value increases, assessment ratios tend to fall. For learning purposes, it is useful to rewrite this statement in a form that is directly compatible with regression outputs.

By \eqref{eq:log_ratio_identity}, the log assessment ratio is exactly the log-scale residual. Hence, if we let
\[
e_i := \log r_i
\qquad\text{and}\qquad
y_i := \log p_i,
\]
then any systematic price-related pattern in the ratios can equivalently be studied as a relationship between residuals and log prices. A convenient sample summary of this relationship on an evaluation sample $S$ is
\begin{equation}
\label{eq:sample_cov}
\widehat{\operatorname{Cov}}_S(e,y)
:=
\frac{1}{|S|}
\sum_{i\in S}
\big(e_i-\bar e_S\big)\big(y_i-\bar y_S\big),
\end{equation}
where $\bar e_S$ and $\bar y_S$ denote the sample means of $\{e_i\}_{i\in S}$ and $\{y_i\}_{i\in S}$, respectively. The corresponding slope of log ratios on log prices is
\begin{equation}
\label{eq:sample_slope}
\widehat{\beta}^{\,\log}_S
:=
\frac{\widehat{\operatorname{Cov}}_S(e,y)}
{\widehat{\operatorname{Var}}_S(y)}.
\end{equation}
Because $\widehat{\operatorname{Var}}_S(y)$ is fixed once the evaluation sample is fixed, controlling the sign and magnitude of price-related linear trends is closely tied to controlling the covariance between log ratios and log prices.

This training-oriented summary is \emph{not} intended to replace assessor metrics such as PRD or PRB. Rather, it serves a different role: it provides a smooth, model-compatible summary of price-related dependence that can be acted on during learning. In this sense, the paper moves from \emph{diagnosing} regressivity through ratio-study statistics to \emph{controlling} a closely related dependence pattern during training.

\medskip
\noindent\textbf{A population intuition under squared-error prediction.}
The dependence view is also consistent with a basic population property of squared-error regression. Let $P$ denote the random market value, let $Y:=\log P$, and let $X$ denote the observed feature vector. Under squared loss, the Bayes predictor is the conditional mean
\[
f^\star(X)=\mathbb{E}[Y\mid X].
\]
If we define the associated residual by
\[
e^\star := f^\star(X)-Y,
\]
then
\begin{equation}
\label{eq:population_identity}
\operatorname{Cov}(e^\star,Y)
=
-\mathbb{E}\!\left[\operatorname{Var}(Y\mid X)\right]
\le 0.
\end{equation}
Thus, whenever some irreducible conditional variation in prices remains, the conditional-mean predictor induces a non-positive association between residuals and log prices. Since $e^\star=\log(\hat P^\star/P)$, the same statement can be read as a non-positive association between log assessment ratios and log values.

This identity should be interpreted as an intuition, not as a claim that every finite-sample model trained for predictive accuracy must be regressive under every assessor metric. Its role is to explain why accuracy-oriented regression naturally tends to compress the extremes of the value distribution, especially when important price drivers are unobserved or imperfectly measured. In mass appraisal, this is exactly the mechanism that can translate into overvaluation at the low end and undervaluation at the high end.

\subsection{Temporal Evaluation and Validation Protocol}
\label{subsec:oos_protocol}

Mass appraisal is a temporally deployed prediction problem: models are trained on historical sales and then used to support later assessments. For this reason, the relevant notion of out-of-sample performance is chronological rather than purely random-split performance. IAAO explicitly recommends holdout-based model testing for objective performance assessment, and rolling-origin evaluation is standard in the forecasting literature when model development must respect time order \cite{IAAO2025Mass,tashman2000out}.

Let the full sold-data sample be ordered by sale date and partitioned into two chronologically separated parts:
\[
\mathcal{D}^{\mathrm{dev}}
\qquad\text{and}\qquad
\mathcal{D}^{\mathrm{test}},
\]
where $\mathcal{D}^{\mathrm{dev}}$ is the development sample used for training and calibration, and $\mathcal{D}^{\mathrm{test}}$ is a final holdout assessment period used only for one-time reporting after model selection is complete. No model-selection decision is allowed to use information from $\mathcal{D}^{\mathrm{test}}$.

Within the development sample, we construct a sequence of rolling-origin windows
\[
\big\{(\mathcal{T}_w,\mathcal{V}_w)\big\}_{w=1}^W,
\]
where each training set $\mathcal{T}_w$ precedes its corresponding validation set $\mathcal{V}_w$ in time. The windows may be expanding or otherwise office-specific, but they always respect the chronological requirement
\[
\max_{i\in \mathcal{T}_w} t_i
<
\min_{j\in \mathcal{V}_w} t_j.
\]
A \emph{candidate configuration} $c\in\mathcal{C}$ consists of a model class together with its tuning choices, including any regressivity-aware penalty parameters introduced later. For each $c\in\mathcal{C}$ and each window $w$, we fit the model on $\mathcal{T}_w$, predict on $\mathcal{V}_w$, and record the corresponding out-of-sample performance.

To keep the notation flexible, let
\[
\mathbf{A}_w(c)
\quad\text{and}\quad
\mathbf{Q}_w(c)
\]
denote the vectors of accuracy and equity summaries, respectively, computed on the validation sample $\mathcal{V}_w$. Accuracy metrics are evaluated on the prediction scale specified in the empirical section. Equity metrics are always computed after back-transforming predictions to the original price scale and forming the ratios in \eqref{eq:ratio_def}. This distinction is important and will be maintained throughout the paper: models are trained on log prices, but assessor-facing fairness is judged on price-scale ratios.

In this paper, \emph{out-of-sample} means evaluation on sales whose dates are strictly later than the dates of the sales used to fit the corresponding model. Section~\ref{sec:model_selection} uses the collection
\[
\Big\{ \mathbf{A}_w(c),\mathbf{Q}_w(c)\Big\}_{c\in\mathcal{C},\,w=1,\dots,W}
\]
to define the multi-objective calibration and model-selection rule. After that rule selects a final configuration, the model is refit on the full development sample $\mathcal{D}^{\mathrm{dev}}$ and evaluated once on the untouched holdout sample $\mathcal{D}^{\mathrm{test}}$.

The purpose of the present section is therefore twofold. First, it defines the ratio-study objects through which regressivity and assessment equity are measured. Second, it defines the temporal evaluation environment in which competing models will be compared. Section~\ref{sec:methodology} then uses this setup to introduce regressivity-aware regularization terms for standard regression learners.








\section{General Penalized Learning Framework}

We now formalize the in-processing approach introduced in the previous section. The goal is to augment a standard predictive learner with an additional term that discourages systematic value-dependent distortions in the resulting assessments, while preserving the basic structure of the underlying regression model.

Let $\D := \{(x_i,y_i)\}_{i=1}^n$ be a training sample, where $x_i \in \X$ denotes the feature vector of property $i$ and $y_i \in \R$ denotes its log market value. Let $\H$ be a hypothesis class of prediction functions $f:\X \to \R$. The baseline predictive model is defined by the empirical risk minimization problem
\begin{align}
    \label{opt:base_model}
    f_{\mathrm{base}}
    \in
    \arg\min_{f \in \H}
    \left\{
        L_n(f) + \mathcal{R}(f)
    \right\},
    \qquad
    L_n(f) := \frac{1}{n}\sum_{i=1}^n \ell(f(x_i),y_i),
\end{align}
where $\ell$ is a prescribed loss function and $\mathcal{R}:\H \to \R$ is an optional complexity regularizer. Throughout this section, all quantities are empirical unless otherwise stated.

To connect this predictive problem with assessment equity, define the fitted log value
\[
\hat y_i(f) := f(x_i),
\]
and the log-scale residual
\[
e_i(f) := \hat y_i(f) - y_i.
\]
Since $y_i = \log P_i$ and $\hat y_i(f) = \log \hat P_i(f)$, where $P_i$ and $\hat P_i(f)$ denote the market value and fitted assessed value of property $i$, respectively, we have
\begin{align}
    \label{eq:log_ratio_identity}
    e_i(f)
    = \log \hat P_i(f) - \log P_i
    = \log\!\left(\frac{\hat P_i(f)}{P_i}\right)
    = \log r_i(f),
\end{align}
where
\[
r_i(f) := \frac{\hat P_i(f)}{P_i}
\]
is the assessment-to-sale ratio induced by the predictor $f$. Thus, controlling systematic dependence between the residuals $e_i(f)$ and the log values $y_i$ is equivalent to controlling systematic dependence between log assessment ratios and log market values.

Motivated by this equivalence, let $\Psi_n(f)$ denote a sample-dependent functional designed to measure regressivity-relevant dependence in the fitted assessments. One may then formulate a constrained learning problem of the form
\begin{align}
    \label{opt:general_constrained_model}
    \min_{f \in \H}
    \left\{
        L_n(f) + \mathcal{R}(f)
    \right\}
    \quad
    \text{s.t.}
    \quad
    \Psi_n(f) \le \tau,
\end{align}
for some tolerance level $\tau \ge 0$. In this paper, we work instead with the corresponding soft-constrained penalized formulation
\begin{align}
    \label{opt:general_penalized_model}
    f_{\rho}
    \in
    \arg\min_{f \in \H}
    \left\{
        L_n(f) + \mathcal{R}(f) + \rho\,\Psi_n(f)
    \right\},
\end{align}
where $\rho \ge 0$ is a penalty parameter that controls the strength of the regressivity correction. The baseline model \eqref{opt:base_model} is recovered as the special case $\rho=0$.

Formulation \eqref{opt:general_penalized_model} isolates the general structure of the proposed approach: starting from a standard predictive learner, we augment the empirical objective with a penalty that discourages systematic value-dependent distortions in the resulting assessments. In the present setting, the relevant distortions are naturally expressed through the dependence between assessment ratios and property values, or equivalently, between log-ratios and log-prices. The next section makes this idea concrete by introducing covariance-based choices of the functional $\Psi_n$, including both direct covariance penalties and computationally convenient surrogate forms.

\paragraph{Remark.}
Once a concrete choice of $\Psi_n$ is specified, problem \eqref{opt:general_penalized_model} can be instantiated in standard learners used in mass appraisal, including regularized linear models and tree-based boosting methods. These model-specific implementations are developed after the covariance-based penalties have been introduced.

\section{Covariance-Based Penalizations}

We now instantiate the general penalized learning problem \eqref{opt:general_penalized_model} by specifying concrete choices of the functional $\Psi_n(f)$ based on empirical covariance-type measures. Throughout, we keep the notation of the previous section: for a predictor $f \in \H$,
\[
\hat y_i(f) = f(x_i),
\qquad
e_i(f) = \hat y_i(f) - y_i,
\qquad
r_i(f) = \exp(e_i(f)),
\]
so that $e_i(f)=\log r_i(f)$ by \eqref{eq:log_ratio_identity}. Since regressivity corresponds to a systematic decrease of assessment ratios with property values, a natural strategy is to penalize dependence between the vectors $e(f):=(e_1(f),\dots,e_n(f))$ and $y:=(y_1,\dots,y_n)$.

Let
\[
\bar e(f) := \frac{1}{n}\sum_{i=1}^n e_i(f),
\qquad
\bar y := \frac{1}{n}\sum_{i=1}^n y_i,
\]
and define the empirical covariance
\begin{align}
    \label{eq:empirical_covariance}
    \widehat{\Cov}_n(e(f),y)
    :=
    \frac{1}{n}\sum_{i=1}^n
    \bigl(e_i(f)-\bar e(f)\bigr)\bigl(y_i-\bar y\bigr).
\end{align}
This quantity is directly connected to the slope of the least-squares regression of $e_i(f)$ on $y_i$, since
\[
\widehat{\beta}_{e\sim y}(f)
=
\frac{\widehat{\Cov}_n(e(f),y)}{\widehat{\Var}_n(y)},
\]
and $\widehat{\Var}_n(y)$ does not depend on the model choice. Accordingly, controlling the empirical covariance is a natural way to control first-order value-dependent trends in the fitted assessment ratios.

\subsection{Direct covariance penalties}

A first class of penalties is obtained by using the empirical covariance in \eqref{eq:empirical_covariance} directly.

\paragraph{Squared covariance penalty.}
The most symmetric choice is to penalize the squared covariance:
\begin{align}
    \label{eq:cov2_penalty}
    \Psi_n^{\mathrm{cov}^2}(f)
    :=
    \frac{1}{2}\,
    \widehat{\Cov}_n(e(f),y)^2.
\end{align}
Substituting \eqref{eq:cov2_penalty} into \eqref{opt:general_penalized_model} yields
\begin{align}
    \label{opt:cov2_penalized_model}
    f_{\rho}^{\mathrm{cov}^2}
    \in
    \arg\min_{f\in\H}
    \left\{
        L_n(f)+\mathcal{R}(f)
        +
        \frac{\rho}{2}\,\widehat{\Cov}_n(e(f),y)^2
    \right\}.
\end{align}
This choice penalizes any linear dependence between $e(f)$ and $y$, regardless of sign. In particular, it suppresses both regressive and progressive linear trends, and is therefore appropriate when the goal is to drive the first-order dependence toward zero.

\paragraph{Signed covariance penalty.}
In some applications, one may wish to focus specifically on regressive patterns under the sign convention adopted here, namely $e_i(f)=\log r_i(f)$. Since regressivity corresponds to negative dependence between $e(f)$ and $y$, a directional alternative is
\begin{align}
    \label{eq:negcov_penalty}
    \Psi_n^{\mathrm{neg\mbox{-}cov}}(f)
    :=
    -\,\widehat{\Cov}_n(e(f),y).
\end{align}
The resulting optimization problem is
\begin{align}
    \label{opt:negcov_penalized_model}
    f_{\rho}^{\mathrm{neg\mbox{-}cov}}
    \in
    \arg\min_{f\in\H}
    \left\{
        L_n(f)+\mathcal{R}(f)
        -
        \rho\,\widehat{\Cov}_n(e(f),y)
    \right\}.
\end{align}
Unlike \eqref{opt:cov2_penalized_model}, this formulation is directional rather than symmetric: it directly pushes the fitted model away from negative covariance and toward less regressive dependence structures. We therefore view \eqref{opt:negcov_penalized_model} as a sign-sensitive alternative when only regressive trends, rather than all linear trends, are the primary concern.

\subsection{A separable surrogate}

While the direct covariance penalties above are conceptually natural, they have two important limitations. First, covariance is a first-order summary and may fail to reflect more complex dependence patterns when positive and negative local trends cancel each other out. Second, the term $\widehat{\Cov}_n(e(f),y)^2$ is non-separable across observations, which can be inconvenient in optimization methods that rely on observation-wise gradients or second-order approximations.

To motivate a more computationally convenient alternative, define
\[
a_i(f):=\bigl(e_i(f)-\bar e(f)\bigr)\bigl(y_i-\bar y\bigr).
\]
Then,
\[
\widehat{\Cov}_n(e(f),y)
=
\frac{1}{n}\sum_{i=1}^n a_i(f),
\]
and Jensen's inequality implies
\begin{align}
    \label{eq:covariance_upper_bound}
    \widehat{\Cov}_n(e(f),y)^2
    =
    \left(
        \frac{1}{n}\sum_{i=1}^n a_i(f)
    \right)^2
    \le
    \frac{1}{n}\sum_{i=1}^n a_i(f)^2
    =
    \frac{1}{n}\sum_{i=1}^n
    \bigl(e_i(f)-\bar e(f)\bigr)^2
    \bigl(y_i-\bar y\bigr)^2.
\end{align}
The right-hand side of \eqref{eq:covariance_upper_bound} is an upper bound on the squared covariance, but it still depends on the sample mean $\bar e(f)$ and therefore remains coupled across observations.

In many well-calibrated regression settings, however, the average residual is expected to be small, and perfect equity corresponds to $e_i(f)=0$ for every $i$, or equivalently $r_i(f)=1$. Motivated by this observation, we consider the simplified surrogate
\begin{align}
    \label{eq:surrogate_penalty}
    \Psi_n^{\mathrm{surr}}(f)
    :=
    \frac{1}{n}\sum_{i=1}^n
    e_i(f)^2\bigl(y_i-\bar y\bigr)^2.
\end{align}
Substituting \eqref{eq:surrogate_penalty} into \eqref{opt:general_penalized_model} gives the surrogate-penalized problem
\begin{align}
    \label{opt:surrogate_penalized_model}
    f_{\rho}^{\mathrm{surr}}
    \in
    \arg\min_{f\in\H}
    \left\{
        L_n(f)+\mathcal{R}(f)
        +
        \rho\,\Psi_n^{\mathrm{surr}}(f)
    \right\}.
\end{align}

The surrogate \eqref{eq:surrogate_penalty} has two useful features. First, it remains anchored to the same underlying idea as the covariance penalties: residual distortions are weighted more heavily when they occur far from the center of the value distribution. Second, unlike \eqref{eq:cov2_penalty}, it is separable across observations once the weights $(y_i-\bar y)^2$ are fixed. This makes \eqref{opt:surrogate_penalized_model} particularly attractive for model classes and optimization procedures in which observation-wise objective contributions are computationally convenient.

\subsection{Discussion}

The three formulations above target related but distinct notions of value-dependent distortion. The squared covariance penalty \eqref{eq:cov2_penalty} is the most direct way to suppress linear dependence between log-ratios and log-values. The signed covariance penalty \eqref{eq:negcov_penalty} focuses instead on the direction of dependence associated with regressivity under the present sign convention. The surrogate penalty \eqref{eq:surrogate_penalty} sacrifices exact covariance control in exchange for a data-point-wise separable objective that is better suited to some learning algorithms.

In the remainder of the paper, these penalties are instantiated within standard regression models used in practice by assessors. In particular, the next section shows how the covariance-based objectives above can be implemented in model classes such as regularized linear regression and gradient boosting, thereby preserving the general learning framework while yielding practically usable regressivity-aware predictors.

\section{Model-Specific Implementations}
\label{sec:model_specific_implementations}

We now show how the covariance-based penalizations introduced in the previous section can be instantiated in specific model classes. We focus on two examples that are especially relevant in practice: linear models and second-order tree boosting. The former provides a transparent analytical benchmark, while the latter covers widely used high-performing nonlinear learners such as XGBoost- and LightGBM-type methods.

\subsection{Linear models}
\label{subsec:linear_models}

Consider the linear hypothesis class
\[
\H_{\mathrm{lin}} := \{f_\theta : x \mapsto x^\top \theta \;:\; \theta \in \Theta \subseteq \R^p\},
\]
and let $X \in \R^{n\times p}$ denote the design matrix. For a coefficient vector $\theta \in \Theta$, define
\[
f_\theta(x_i) = x_i^\top \theta,
\qquad
e(\theta) := X\theta - y.
\]
As in the previous section, let
\[
H := I - \frac{1}{n}\mathbf{1}\mathbf{1}^\top,
\qquad
c := Hy = y - \bar y\,\mathbf{1},
\qquad
D := \diag(c_1^2,\dots,c_n^2),
\]
so that
\[
\widehat{\Cov}_n(e(\theta),y)
=
\frac{1}{n}e(\theta)^\top c.
\]

Throughout this subsection, we take the baseline loss to be the squared loss
\[
L_n(f_\theta) = \frac{1}{2n}\|X\theta-y\|_2^2.
\]
We then derive the corresponding penalized estimators for the three covariance-based choices introduced in Section~\ref{sec:covariance_penalizations}.

\paragraph{Signed covariance penalty.}
For the directional penalty
\[
\Psi_n^{\mathrm{neg\mbox{-}cov}}(f_\theta)
=
-\widehat{\Cov}_n(e(\theta),y)
=
-\frac{1}{n}e(\theta)^\top c,
\]
the penalized objective is
\begin{align}
    \label{opt:linear_negcov}
    \theta_{\rho}^{\mathrm{neg\mbox{-}cov}}
    \in
    \arg\min_{\theta\in\Theta}
    \left\{
        \frac{1}{2n}\|X\theta-y\|_2^2
        -
        \frac{\rho}{n}e(\theta)^\top c
    \right\}.
\end{align}
Completing the square gives
\[
\frac{1}{2n}\|X\theta-y\|_2^2
-
\frac{\rho}{n}e(\theta)^\top c
=
\frac{1}{2n}\|X\theta-(y+\rho c)\|_2^2
-
\frac{\rho^2}{2n}\|c\|_2^2,
\]
so \eqref{opt:linear_negcov} is equivalent, up to an additive constant, to ordinary least squares with shifted response $y+\rho c$. Consequently, whenever $X^\top X$ is invertible,
\begin{align}
    \label{eq:linear_negcov_closed_form}
    \theta_{\rho}^{\mathrm{neg\mbox{-}cov}}
    =
    (X^\top X)^{-1}X^\top (y+\rho c).
\end{align}
Thus, under squared loss, the signed covariance penalty is particularly simple to implement in linear models.

\paragraph{Squared covariance penalty.}
For the symmetric penalty
\[
\Psi_n^{\mathrm{cov}^2}(f_\theta)
=
\frac{1}{2}\,\widehat{\Cov}_n(e(\theta),y)^2
=
\frac{1}{2n^2}\bigl(e(\theta)^\top c\bigr)^2,
\]
the penalized estimator solves
\begin{align}
    \label{opt:linear_cov2}
    \theta_{\rho}^{\mathrm{cov}^2}
    \in
    \arg\min_{\theta\in\Theta}
    \left\{
        \frac{1}{2n}\|X\theta-y\|_2^2
        +
        \frac{\rho}{2n^2}\bigl((X\theta-y)^\top c\bigr)^2
    \right\}.
\end{align}
Define
\[
a := X^\top c,
\qquad
b := y^\top c = c^\top c.
\]
Then the first-order optimality condition for \eqref{opt:linear_cov2} is
\[
\left(
\frac{1}{n}X^\top X + \frac{\rho}{n^2}aa^\top
\right)\theta
=
\frac{1}{n}X^\top y + \frac{\rho}{n^2}ba.
\]
Hence, whenever the displayed inverse exists,
\begin{align}
    \label{eq:linear_cov2_closed_form}
    \theta_{\rho}^{\mathrm{cov}^2}
    =
    \left(
        \frac{1}{n}X^\top X
        +
        \frac{\rho}{n^2}aa^\top
    \right)^{-1}
    \left(
        \frac{1}{n}X^\top y
        +
        \frac{\rho}{n^2}ba
    \right).
\end{align}
This shows that the squared covariance penalty induces a rank-one modification of the usual normal equations.

\paragraph{Separable surrogate penalty.}
For the separable surrogate
\[
\Psi_n^{\mathrm{surr}}(f_\theta)
=
\frac{1}{n}\sum_{i=1}^n e_i(\theta)^2 c_i^2
=
\frac{1}{n}e(\theta)^\top D\,e(\theta),
\]
the penalized problem is
\begin{align}
    \label{opt:linear_surr}
    \theta_{\rho}^{\mathrm{surr}}
    \in
    \arg\min_{\theta\in\Theta}
    \left\{
        \frac{1}{2n}\|X\theta-y\|_2^2
        +
        \frac{\rho}{n}e(\theta)^\top D\,e(\theta)
    \right\}.
\end{align}
The corresponding optimality condition is
\[
\left(
\frac{1}{n}X^\top X
+
\frac{2\rho}{n}X^\top D X
\right)\theta
=
\frac{1}{n}X^\top y
+
\frac{2\rho}{n}X^\top D y,
\]
and therefore, whenever the displayed inverse exists,
\begin{align}
    \label{eq:linear_surr_closed_form}
    \theta_{\rho}^{\mathrm{surr}}
    =
    \left(
        \frac{1}{n}X^\top X
        +
        \frac{2\rho}{n}X^\top D X
    \right)^{-1}
    \left(
        \frac{1}{n}X^\top y
        +
        \frac{2\rho}{n}X^\top D y
    \right).
\end{align}
The factor \(2\rho\) arises directly from the definition of \(\Psi_n^{\mathrm{surr}}\) in Section~\ref{sec:covariance_penalizations}.

\paragraph{Extension to regularized linear models.}
The same constructions extend immediately to penalized linear estimators of the form
\[
\min_{\theta\in\Theta}
\left\{
    L_n(f_\theta) + \mathcal{R}(\theta) + \rho\,\Psi_n(f_\theta)
\right\},
\]
where \(\mathcal{R}\) is a complexity regularizer on the coefficients. In the ridge case, the resulting systems remain quadratic and admit closed-form solutions after adding the usual ridge term to the coefficient matrix. In contrast, \(\ell_1\)-regularized variants such as the LASSO remain convex but no longer admit closed-form expressions, and must instead be solved numerically.

\subsection{Second-order Gradient boosting}
\label{subsec:tree_boosting}

We next consider second-order additive tree boosting under squared loss, as used in gradient-boosted tree systems such as XGBoost and LightGBM. Let \(f^{(m-1)}\) denote the current prediction function at boosting iteration \(m\), and write
\[
f_i := f^{(m-1)}(x_i),
\qquad
e_i := f_i - y_i,
\qquad
c_i := y_i - \bar y,
\qquad
z := \widehat{\Cov}_n(e,y) = \frac{1}{n}\sum_{i=1}^n e_i c_i.
\]
At step \(m\), one adds a regression tree \(t\) to the current predictor, so that
\[
f^{(m)}(x) = f^{(m-1)}(x) + \nu\,t(x),
\]
where \(\nu\in(0,1]\) is the learning rate.

For standard second-order boosting, the objective is approximated locally by
\begin{align}
    \label{eq:boosting_local_objective}
    \widetilde{\mathcal J}(t)
    =
    \sum_{i=1}^n
    \left(
        g_i\,t_i + \frac{1}{2}h_i\,t_i^2
    \right)
    +
    \Omega(t),
\end{align}
where \(t_i := t(x_i)\), \(\Omega(t)\) is the tree-complexity penalty, and
\[
g_i := \frac{\partial}{\partial f_i}
\left(
L_n(f) + \rho\,\Psi_n(f)
\right),
\qquad
h_i := \frac{\partial^2}{\partial f_i^2}
\left(
L_n(f) + \rho\,\Psi_n(f)
\right)
\]
are evaluated at the current predictions. Under squared loss,
\[
\frac{\partial L_n(f)}{\partial f_i} = \frac{e_i}{n},
\qquad
\frac{\partial^2 L_n(f)}{\partial f_i^2} = \frac{1}{n}.
\]
The covariance-based penalties modify these quantities as follows.

\paragraph{Signed covariance penalty.}
For \(\Psi_n^{\mathrm{neg\mbox{-}cov}}(f) = -n^{-1}\sum_{i=1}^n e_i c_i\), we obtain
\begin{align}
    \label{eq:boosting_negcov_grad_hess}
    g_i^{\mathrm{neg\mbox{-}cov}}
    =
    \frac{e_i}{n}
    -
    \frac{\rho}{n}c_i,
    \qquad
    h_i^{\mathrm{neg\mbox{-}cov}}
    =
    \frac{1}{n}.
\end{align}
Thus, the signed covariance penalty yields an exact observation-wise objective and can be implemented directly through standard custom-gradient interfaces.

\paragraph{Squared covariance penalty.}
For \(\Psi_n^{\mathrm{cov}^2}(f)=\tfrac12 z^2\), the gradient is
\begin{align}
    \label{eq:boosting_cov2_gradient}
    g_i^{\mathrm{cov}^2}
    =
    \frac{e_i}{n}
    +
    \frac{\rho}{n}z\,c_i.
\end{align}
The exact Hessian of the penalized objective with respect to the prediction vector \(f=(f_1,\dots,f_n)\) is
\begin{align}
    \label{eq:boosting_cov2_exact_hessian}
    \nabla_f^2
    \bigl(L_n(f)+\rho\,\Psi_n^{\mathrm{cov}^2}(f)\bigr)
    =
    \frac{1}{n}I
    +
    \frac{\rho}{n^2}cc^\top.
\end{align}
Unlike the separable boosting objectives used in standard libraries, \eqref{eq:boosting_cov2_exact_hessian} is dense: each observation is coupled to every other observation through the rank-one term \(cc^\top\). Consequently, there are two natural implementation routes.

The first is a practical approximation that keeps only the diagonal terms,
\begin{align}
    \label{eq:boosting_cov2_diag_hessian}
    h_i^{\mathrm{cov}^2,\mathrm{diag}}
    =
    \frac{1}{n}
    +
    \frac{\rho}{n^2}c_i^2.
\end{align}
This produces a custom objective compatible with standard second-order boosting APIs, at the price of ignoring the cross-observation curvature induced by the covariance penalty.

The second is an exact leaf-wise implementation for a fixed tree structure. Suppose a candidate tree \(t\) has leaves \(I_1,\dots,I_J\) and leaf weights \(w=(w_1,\dots,w_J)\), so that \(t_i=w_j\) for \(i\in I_j\). Let
\[
n_j := |I_j|,
\qquad
b_j := \sum_{i\in I_j} e_i,
\qquad
s_j := \sum_{i\in I_j} c_i,
\]
and let \(s=(s_1,\dots,s_J)^\top\). If \(\Omega(t)=\gamma J + \frac{\lambda}{2}\|w\|_2^2\), then the exact local objective for the squared covariance penalty can be written as
\begin{align}
    \label{eq:boosting_cov2_exact_leaf_objective}
    \widetilde{\mathcal J}_{\mathrm{cov}^2}(w)
    =
    g_T^\top w
    +
    \frac{1}{2}w^\top A_T w
    +
    \gamma J
    + \text{const},
\end{align}
where
\[
g_{T,j}
=
\frac{b_j}{n}
+
\frac{\rho}{n}z\,s_j,
\qquad
A_T
=
\diag\!\left(\frac{n_1}{n}+\lambda,\dots,\frac{n_J}{n}+\lambda\right)
+
\frac{\rho}{n^2}ss^\top.
\]
Hence the optimal leaf weights satisfy
\begin{align}
    \label{eq:boosting_cov2_exact_leaf_system}
    A_T w^\star = -g_T,
    \qquad
    w^\star = -A_T^{-1}g_T,
\end{align}
whenever \(A_T\) is invertible. The associated gain for the fixed structure \(T\) is
\begin{align}
    \label{eq:boosting_cov2_exact_gain}
    \operatorname{Gain}(T)
    =
    \frac{1}{2}g_T^\top A_T^{-1}g_T - \gamma J.
\end{align}
This exact route preserves the full curvature of the covariance penalty, but requires a custom tree-construction procedure rather than the usual observation-wise Hessian interface.

\paragraph{Separable surrogate penalty.}
For the surrogate
\[
\Psi_n^{\mathrm{surr}}(f)
=
\frac{1}{n}\sum_{i=1}^n e_i^2 c_i^2,
\]
the gradient and Hessian remain observation-wise separable:
\begin{align}
    \label{eq:boosting_surr_grad_hess}
    g_i^{\mathrm{surr}}
    &=
    \frac{e_i}{n}
    +
    \frac{2\rho}{n}e_i c_i^2,
    \\
    h_i^{\mathrm{surr}}
    &=
    \frac{1}{n}
    +
    \frac{2\rho}{n}c_i^2.
\end{align}
Accordingly, the surrogate penalty is particularly well suited to off-the-shelf second-order boosting implementations: once the weights \(c_i^2\) are precomputed, the penalized objective can be supplied through standard custom gradient and Hessian routines.

\paragraph{Implementation summary.}
The three penalties above lead to distinct computational regimes in boosted trees. The signed covariance penalty is fully separable and is therefore the simplest directional correction to implement. The surrogate penalty is also separable, while placing greater emphasis on distortions occurring away from the center of the value distribution. The squared covariance penalty is the most direct implementation of the covariance-control idea, but introduces dense curvature; in practice, it can be handled either through a diagonal Hessian approximation in standard boosting libraries or through an exact custom leaf solver of the form \eqref{eq:boosting_cov2_exact_leaf_system}.

The formulas in this section show that the covariance-based framework is compatible with both transparent linear estimators and high-performing nonlinear learners. This model-specific flexibility is important for the practical goal of the paper: regressivity correction can be layered onto standard predictive tools already used in assessment workflows, rather than requiring a separate valuation architecture.

\paragraph{Algorithmic details.}
For completeness, Appendix~\ref{app:boosting_algorithms} provides pseudocode for an exact Newton-style tree-boosting implementation of the squared covariance penalty, together with an optional first-order variant. In the main text, we focus on the gradient, Hessian, and leaf-system formulas that define the method and clarify its relation to standard boosted-tree implementations.

\section{Constrained Model Selection Across Validation Folds}
\label{subsec:model_selection}

The covariance-regularized learners developed in the previous sections enlarge the set of candidate specifications available at calibration time, but they do not alter the basic structure of the assessor's workflow. After fitting the implemented model families---namely regularized linear models and gradient boosting models, with or without covariance-based penalization---the practical task remains to compare a finite collection of candidate specifications on out-of-sample validation data and to select one final specification for deployment. The difference is that, once regressivity-aware penalties are introduced, candidate models may trade off predictive performance, ratio-study diagnostics, regressivity measures, and dispersion measures in different ways. Accordingly, model calibration is naturally formulated as a constrained multi-criteria selection problem.

\paragraph{Single-model calibration.}
Let $\mathcal V=\{1,\dots,V\}$ index the validation folds (or rolling-origin validation windows), and let $V_v$ denote the validation sample associated with fold $v\in\mathcal V$. Let $\mathcal C$ denote a finite set of candidate model configurations. Each configuration $c\in\mathcal C$ specifies
(i) a model family from the implemented classes,
(ii) a penalty choice
\[
\Psi_n \in \{0,\Psi_n^{\mathrm{cov}^2},\Psi_n^{\mathrm{neg\mbox{-}cov}},\Psi_n^{\mathrm{surr}}\},
\]
where $\Psi_n=0$ denotes the unpenalized baseline,
and
(iii) all associated tuning parameters, including the penalty level $\rho\ge 0$ when applicable.
For each fold $v\in\mathcal V$, fitting configuration $c$ on the corresponding training portion yields a predictor $f_{c,v}$.

For each observation $i\in V_v$, let
\[
\hat y_i^{\,v}(c):=f_{c,v}(x_i^{\,v}),
\qquad
\hat P_i^{\,v}(c):=\exp\!\bigl(\hat y_i^{\,v}(c)\bigr),
\qquad
r_i^{\,v}(c):=\frac{\hat P_i^{\,v}(c)}{P_i^{\,v}},
\]
where $P_i^{\,v}$ is the observed sale price. From the validation predictions on fold $v$, we compute a predictive performance score $\mathrm{Perf}_v(c)$, written so that larger values are preferred, together with a collection of fold-level assessment diagnostics $Q_{m,v}(c)$ indexed by $m$. If one starts from an error metric $E_v(c)$ that is minimized, such as MSE, RMSE, or MAE, one may simply set $\mathrm{Perf}_v(c):=-E_v(c)$.

To aggregate results across folds, let $\alpha_v\ge 0$ satisfy $\sum_{v\in\mathcal V}\alpha_v=1$, and define
\[
\bar{\mathrm{Perf}}(c):=\sum_{v\in\mathcal V}\alpha_v\,\mathrm{Perf}_v(c),
\qquad
\bar Q_m(c):=\sum_{v\in\mathcal V}\alpha_v\,Q_{m,v}(c).
\]
The default choice is uniform weighting, $\alpha_v=1/V$, although the same formulation also allows non-uniform weights when one wishes to emphasize, for example, more recent validation windows.

We then select a single deployed model by solving
\begin{align}
    \label{opt:model_selection_single}
    c^\star
    \in
    \arg\max_{c\in\mathcal C}
    \quad
    & \bar{\mathrm{Perf}}(c)
    \\
    \text{s.t.}\quad
    & \underline{\kappa}_m \le \bar Q_m(c) \le \overline{\kappa}_m,
    \qquad m\in\mathcal M_{\mathrm{int}},
    \nonumber\\
    & \bar Q_m(c) \le \overline{\kappa}_m,
    \qquad m\in\mathcal M_{\mathrm{up}}.
    \nonumber
\end{align}
Here, $\mathcal M_{\mathrm{int}}$ indexes diagnostics constrained to lie in admissible intervals, while $\mathcal M_{\mathrm{up}}$ indexes diagnostics handled through upper bounds only. Typical elements of $\mathcal M_{\mathrm{int}}$ include ratio statistics such as the mean ratio, weighted mean ratio, and median ratio, as well as regressivity diagnostics such as VEI, PRD, and PRB. Typical elements of $\mathcal M_{\mathrm{up}}$ include dispersion measures such as COD and COV. Additional assessor-facing diagnostics can be incorporated by adding them to either set according to the desired feasibility requirement.

Formulation \eqref{opt:model_selection_single} deliberately avoids collapsing all validation criteria into a single weighted score. Predictive performance remains the objective, while ratio-study, regressivity, and dispersion diagnostics enter as explicit feasibility requirements. This mirrors the logic of assessor calibration more closely than a purely accuracy-driven ranking rule: among all candidate specifications that satisfy the chosen assessment-quality requirements, select the one with the best aggregated out-of-sample predictive performance.

\paragraph{How the constrained calibration differs from standard accuracy-only selection.}
The single-model procedure can be read as the following practical recipe.
\begin{enumerate}
    \item Specify a finite candidate set $\mathcal C$ of model configurations, including baseline and covariance-penalized variants.
    \item For each $c\in\mathcal C$ and each validation fold $v\in\mathcal V$, fit the model on the corresponding training portion and compute validation predictions.
    \item On each fold, compute the chosen predictive score $\mathrm{Perf}_v(c)$ together with the selected assessment diagnostics $Q_{m,v}(c)$.
    \item Aggregate these fold-level quantities into $\bar{\mathrm{Perf}}(c)$ and $\bar Q_m(c)$.
    \item Discard all configurations that violate any active assessment-quality constraint.
    \item Rank the remaining feasible configurations by aggregated predictive performance and select the best one.
    \item Refit the selected configuration $c^\star$ on the full pre-assessment development sample and deploy that final predictor.
\end{enumerate}
The only substantive departure from the usual accuracy-only workflow lies in Steps 5--6: instead of ranking all candidates solely by predictive error, the constrained version first screens for feasibility with respect to the chosen assessor-facing diagnostics and then ranks only the feasible candidates by predictive performance.

The formulation above uses aggregated fold-level diagnostics as the default choice. When temporal stability is especially important, one may replace some aggregated constraints with foldwise requirements of the form
\[
Q_{m,v}(c)\in[\underline{\kappa}_m,\overline{\kappa}_m]
\qquad
\forall v\in\mathcal V,
\]
or, more generally, use a more conservative fold aggregator. We keep the weighted-average formulation as the main presentation because it yields a transparent calibration problem over a finite candidate set and aligns naturally with standard validation practice.

\subsubsection{Convex stacking as an extension}
\label{subsubsec:convex_stacking}

A natural extension of \eqref{opt:model_selection_single} is to allow convex combinations of candidate predictions rather than selecting only one configuration. This does not change the objective--constraint logic of the calibration step; it changes only the feasible set. In the single-model case, the decision variable is a discrete configuration \(c\in\mathcal C\), and calibration is therefore a constrained sorting problem over finitely many candidates. In the stacking case, the decision variable becomes a nonnegative weight vector on the simplex, and calibration becomes a constrained optimization problem over those weights.

Let $\mathcal C_{\mathrm{stk}}=\{c_1,\dots,c_M\}\subseteq\mathcal C$ denote the set of candidate configurations admitted into the stacking pool. For each fold $v\in\mathcal V$ and observation $i\in V_v$, let
\[
\hat P_{i,j}^{\,v}:=\hat P_i^{\,v}(c_j)
\]
denote the out-of-fold predicted assessed value on the original price scale produced by configuration $c_j$. For a weight vector
\[
w=(w_1,\dots,w_M)\in\Delta_M
:=
\left\{
w\in\mathbb R_+^M : \sum_{j=1}^M w_j = 1
\right\},
\]
define the stacked prediction on fold $v$ by
\begin{align}
    \label{eq:stacked_prediction}
    \hat P_i^{\,v}(w)
    :=
    \sum_{j=1}^M w_j\,\hat P_{i,j}^{\,v},
    \qquad
    r_i^{\,v}(w)
    :=
    \frac{\hat P_i^{\,v}(w)}{P_i^{\,v}}.
\end{align}
We emphasize that stacking is defined on the \emph{price scale}, not on the log-prediction scale. This choice is natural because the assessment diagnostics entering the constraints are themselves ratio-study quantities defined in terms of assessed values and sale prices.

Given \eqref{eq:stacked_prediction}, one can compute fold-level performance and assessment diagnostics exactly as in the single-model case, now viewed as functions of the weights:
\[
\mathrm{Perf}_v(w),
\qquad
Q_{m,v}(w).
\]
The stacked calibration problem is then
\begin{align}
    \label{opt:model_selection_stacking}
    w^\star
    \in
    \arg\max_{w\in\Delta_M}
    \quad
    & \sum_{v\in\mathcal V}\alpha_v\,\mathrm{Perf}_v(w)
    \\
    \text{s.t.}\quad
    & \underline{\kappa}_m
      \le
      \sum_{v\in\mathcal V}\alpha_v\,Q_{m,v}(w)
      \le
      \overline{\kappa}_m,
      \qquad
      m\in\mathcal M_{\mathrm{int}},
      \nonumber\\
    & \sum_{v\in\mathcal V}\alpha_v\,Q_{m,v}(w)
      \le
      \overline{\kappa}_m,
      \qquad
      m\in\mathcal M_{\mathrm{up}}.
      \nonumber
\end{align}
Thus, the stacking extension preserves the same fold structure, the same diagnostics, and the same feasibility logic as \eqref{opt:model_selection_single}; it simply replaces discrete candidate comparison by continuous optimization over convex combinations of their out-of-fold predictions.

Once an optimal weight vector $w^\star$ has been identified, the final deployed stacked predictor is obtained by refitting each constituent configuration $c_j\in\mathcal C_{\mathrm{stk}}$ on the full pre-assessment development sample and combining their resulting price-scale predictions using the weights $w^\star$.

\subsubsection{Illustrative convex stacking example}
\label{subsubsec:convex_example}

A useful special case is obtained by taking average validation MSE as the objective and imposing worst-fold box constraints on selected ratio-study diagnostics. Let \(n_v:=|V_v|\). For each fold \(v\), define the stacked MSE
\[
\mathrm{MSE}_v(w)
:=
\frac{1}{n_v}\sum_{i\in V_v}
\bigl(\hat P_i^{\,v}(w)-P_i^{\,v}\bigr)^2.
\]
Consider the problem
\begin{align}
    \label{opt:stacking_convex_example}
    \min_{w\in\Delta_M}
    \quad
    & \frac{1}{V}\sum_{v\in\mathcal V}\mathrm{MSE}_v(w)
    \\
    \text{s.t.}\quad
    & \underline{\tau}_{\mathrm{mean}}
      \le
      \frac{1}{n_v}\sum_{i\in V_v}\frac{\hat P_i^{\,v}(w)}{P_i^{\,v}}
      \le
      \overline{\tau}_{\mathrm{mean}},
      \qquad \forall v\in\mathcal V,
      \nonumber\\
    & \underline{\tau}_{\mathrm{wmean}}
      \le
      \frac{\sum_{i\in V_v}\hat P_i^{\,v}(w)}{\sum_{i\in V_v}P_i^{\,v}}
      \le
      \overline{\tau}_{\mathrm{wmean}},
      \qquad \forall v\in\mathcal V,
      \nonumber\\
    & \underline{\tau}_{\mathrm{PRD}}
      \le
      \frac{
        \frac{1}{n_v}\sum_{i\in V_v}\hat P_i^{\,v}(w)/P_i^{\,v}
      }{
        \sum_{i\in V_v}\hat P_i^{\,v}(w)\big/\sum_{i\in V_v}P_i^{\,v}
      }
      \le
      \overline{\tau}_{\mathrm{PRD}},
      \qquad \forall v\in\mathcal V.
      \nonumber
\end{align}
The objective in \eqref{opt:stacking_convex_example} is convex quadratic in \(w\), because each \(\hat P_i^{\,v}(w)\) is affine in \(w\). The mean-ratio and weighted-mean-ratio constraints are linear in \(w\). Moreover, the PRD box constraints can be rearranged into linear inequalities on each fold. Indeed, the lower and upper PRD bounds are equivalent to
\begin{align}
    \label{eq:prd_linearized_bounds}
    \frac{1}{n_v}\sum_{i\in V_v}\frac{\hat P_i^{\,v}(w)}{P_i^{\,v}}
    -
    \underline{\tau}_{\mathrm{PRD}}
    \frac{\sum_{i\in V_v}\hat P_i^{\,v}(w)}{\sum_{i\in V_v}P_i^{\,v}}
    &\ge 0,
    \qquad \forall v\in\mathcal V,
    \\
    \frac{1}{n_v}\sum_{i\in V_v}\frac{\hat P_i^{\,v}(w)}{P_i^{\,v}}
    -
    \overline{\tau}_{\mathrm{PRD}}
    \frac{\sum_{i\in V_v}\hat P_i^{\,v}(w)}{\sum_{i\in V_v}P_i^{\,v}}
    &\le 0,
    \qquad \forall v\in\mathcal V,
    \nonumber
\end{align}
and both sides are affine in \(w\). Consequently, \eqref{opt:stacking_convex_example} is a convex quadratic program over the simplex.

This example illustrates the main appeal of the stacking extension. In the single-model case, calibration sorts a finite set of candidate configurations and selects the best feasible one. In the stacking case, the same assessor-facing diagnostics can be enforced while optimizing continuously over nonnegative combinations of candidate predictors, which can enlarge the attainable accuracy--equity trade-off frontier.

% \paragraph{Remark on COD.}
% Exact COD is not included in \eqref{opt:stacking_convex_example} because it depends on the foldwise median ratio, and that median dependence generally breaks the clean convexity structure above. For this reason, COD fits naturally into the discrete single-model screening problem \eqref{opt:model_selection_single}, whereas a convex stacking formulation is most transparent when built from diagnostics such as mean ratio, weighted mean ratio, and PRD. If one wishes to include a dispersion requirement in the stacking problem while preserving convexity, a convex surrogate or alternative dispersion measure can be used instead.



\section{Constrained Model Selection Across Validation Folds}
\label{sec:model_selection}

The covariance-regularized learners developed in the previous sections enlarge the set of candidate specifications available at calibration time, but they do not change the basic structure of the assessor's workflow. After training the implemented model families---namely regularized linear models and gradient boosting models, with or without covariance-based penalization---the practical task remains to compare a finite collection of candidate specifications on out-of-sample validation data and to select one final specification for deployment. Because the resulting tradeoffs involve predictive performance, ratio-study diagnostics, regressivity measures, and dispersion measures, we formulate this step as a constrained single-model selection problem.

Let $\mathcal V=\{1,\dots,V\}$ index the validation folds (or rolling-origin validation windows), and let $V_v$ denote the validation sample associated with fold $v\in\mathcal V$. Let $\mathcal C$ denote a finite set of candidate model configurations. Each configuration $c\in\mathcal C$ specifies
(i) a model family from the implemented classes,
(ii) a penalty choice
\[
\Psi_n \in \{0,\Psi_n^{\mathrm{cov}^2},\Psi_n^{\mathrm{neg\mbox{-}cov}},\Psi_n^{\mathrm{surr}}\},
\]
where $\Psi_n=0$ denotes the unpenalized baseline, and
(iii) all associated tuning parameters, including the penalty level $\rho\ge 0$ when applicable.
For each fold $v\in\mathcal V$, fitting configuration $c$ on the corresponding training portion yields a predictor $f_{c,v}$.

For each observation $i\in V_v$, let
\[
\hat y_i^{\,v}(c):=f_{c,v}(x_i^{\,v}),
\qquad
\hat P_i^{\,v}(c):=\exp\!\bigl(\hat y_i^{\,v}(c)\bigr),
\qquad
r_i^{\,v}(c):=\frac{\hat P_i^{\,v}(c)}{P_i^{\,v}},
\]
where $P_i^{\,v}$ is the observed sale price. From the validation predictions on fold $v$, we compute a predictive performance score $\mathrm{Perf}_v(c)$, written so that larger values are preferred, together with a collection of fold-level assessment diagnostics $Q_{m,v}(c)$ indexed by $m$. If one starts from an error metric $E_v(c)$ that is minimized, such as MSE, RMSE, or MAE, one may simply set $\mathrm{Perf}_v(c):=-E_v(c)$.

To aggregate results across folds, let $\alpha_v\ge 0$ satisfy $\sum_{v\in\mathcal V}\alpha_v=1$, and define
\[
\bar{\mathrm{Perf}}(c):=\sum_{v\in\mathcal V}\alpha_v\,\mathrm{Perf}_v(c),
\qquad
\bar Q_m(c):=\sum_{v\in\mathcal V}\alpha_v\,Q_{m,v}(c).
\]
The default choice is uniform weighting, $\alpha_v=1/V$, although the same formulation also allows non-uniform weights when one wishes to emphasize, for example, more recent validation windows.

We then select a single deployed model by solving
\begin{align}
    \label{opt:model_selection}
    c^\star
    \in
    \arg\max_{c\in\mathcal C}
    \quad
    & \bar{\mathrm{Perf}}(c)
    \\
    \text{s.t.}\quad
    & \underline{\kappa}_m \le \bar Q_m(c) \le \overline{\kappa}_m,
    \qquad m\in\mathcal M_{\mathrm{int}},
    \nonumber\\
    & \bar Q_m(c) \le \overline{\kappa}_m,
    \qquad m\in\mathcal M_{\mathrm{up}}.
    \nonumber
\end{align}
Here, $\mathcal M_{\mathrm{int}}$ indexes diagnostics constrained to lie in admissible intervals, while $\mathcal M_{\mathrm{up}}$ indexes diagnostics handled through upper bounds only. Typical elements of $\mathcal M_{\mathrm{int}}$ include ratio statistics such as the median ratio, mean ratio, and weighted mean ratio, as well as regressivity diagnostics such as VEI, PRD, and PRB. Typical elements of $\mathcal M_{\mathrm{up}}$ include dispersion measures such as COD and COV. Additional assessor-facing diagnostics can be incorporated by adding them to either set according to the desired feasibility requirement.

Formulation \eqref{opt:model_selection} deliberately avoids collapsing all validation criteria into a single composite score. Predictive performance remains the objective, while ratio-study, regressivity, and dispersion diagnostics enter as explicit constraints. This mirrors the logic of assessor calibration more closely than a weighted-score approach: among all candidate specifications that satisfy the chosen assessment-quality requirements, select the one with the best aggregated out-of-sample predictive performance.

The formulation above uses aggregated fold-level diagnostics as the default choice. When temporal stability is especially important, one may replace some aggregated constraints by foldwise requirements, or use a more conservative aggregation rule. We keep the weighted-average formulation as the main presentation because it yields a transparent calibration problem over a finite candidate set and aligns naturally with standard validation practice.

Once an optimal configuration $c^\star$ has been identified, the final deployed predictor is obtained by refitting that same configuration on the full pre-assessment development sample. In this sense, the procedure remains a single-model selection step: it chooses one specification from the implemented model families and one associated set of tuning parameters, rather than forming a composite predictor.

A natural extension, which we do not pursue here, is to replace the discrete feasible set $\mathcal C$ by convex combinations of candidate prediction vectors. Such a convex-stacking extension would preserve the same objective--constraint logic while potentially enlarging the attainable accuracy--equity trade-off frontier.





\section{Experiments Results}



% OLD TEXT

\newpage 
\section{Background \& Definitions: Title may be even the main topic \& its limitations}

Regressivity, CCAO, IAAO definitions.



\subsubsection{Regressivity in ML-based Mass Appraisal}


\textbf{Mean-squared-error regression induces regressivity in assessments.}
Let $Y \in \mathcal{Y}$ denote the random log-price (the logarithm of the market price), and let $X$ denote the vector of observed features. Under squared-error loss, the Bayes-optimal regressor is
$
f^*(X)=\mathbb{E}[Y\mid X].
$ Define the residual random variable $
e = f^*(X)-Y.
$
Then
\begin{align*}
\Cov(e,Y)
&= \Cov\!\big(\mathbb{E}[Y\mid X]-Y,\,Y\big) 
% &= \Cov\!\big(\mathbb{E}[Y\mid X],Y\big)-\Var(Y) \\
= \Var\!\big(\mathbb{E}[Y\mid X]\big)-\Var(Y) 
= -\,\mathbb{E}\!\left[\Var(Y\mid X)\right]
\;\leq\; 0,
\end{align*}
where the last equality follows from the law of total variance. Therefore, whenever there is nonzero conditional variance in $Y$ given $X$ (i.e., $\mathbb{E}[\Var(Y\mid X)]>0$), we have the negative relationship
$
\Cov(e,Y)<0.
$
Now, since the Assessment-to-Sales ratio is defined as $r=\widehat{P}/P$ (predicted price over market price), then
$
\log r = \log \widehat{P}-\log P = e.
$
Hence, the negative relationship exists for the log-ratios and log-prices as $
\Cov(\log r, Y)=\Cov(e,Y)\leq 0,
$
with strict inequality whenever $\mathbb{E}[\Var(Y\mid X)]>0$. Equivalently, the population linear regression slope of $\log r$ on $Y$ is
\[
\frac{\Cov(\log r,Y)}{\Var(Y)} \leq 0.
\]
This shows that, even under the Bayes-optimal predictor, squared-error prediction naturally induces a non-positive association between log-assessment ratios and log-prices whenever some irreducible conditional variation remains in prices. This is especially relevant in real-estate valuation, where prices typically exhibit substantial idiosyncratic variation not fully explained by the observed covariates (i.e., $\mathbb{E}[\Var(Y\mid X)]>0$).

% \section{Approach: Title of our solution directly can be another alternative}

\section{Soft-constrained Regression Models}

Let $\D:=\{(x_i, y_i)\}_{i=1}^n$ be a given training dataset of $n$ properties, each described by $p$ features (e.g., number of bedrooms, proximity to a bus stop, etc.), where $X\in \R^{n\times p}$ is the feature matrix and $y\in \R^n$ is the target vector of log-prices. Let $\H$ be a hypothesis class associated with a regression learning model (e.g., linear regression, gradient boosting, etc.), consisting of functional mappings of the form $f:\X \rightarrow \Y$, where $\X$ denotes the feature space and $\Y$ the target space. In order to obtain the best-fitting predictor $f^* \in \H$, we solve the complexity-regularized \textit{empirical risk minimization} problem
\begin{align}
    \label{opt:base_model}
    f^*:=\arg\min_{f \in \H} \; \frac{1}{n} \sum_{i=1}^n \ell(f(x_i), y_i) + \mathcal{R}(f),
\end{align}
where $\ell$ is a predefined loss function (e.g., squared error), and $\mathcal{R}:\H \rightarrow \R$ is a complexity-control functional (e.g., an $\ell_1$-type regularization term), which is optional but commonly used in practice. We refer to the solution of \eqref{opt:base_model} as the \textit{base model}.

\medskip

In what follows, we are going to discuss different choices of covariance-based regularization terms that can be considered in order to correct linear trends between $e$ and $y$. Recall from Section/Equation (\todo{refer properly to the cov/var for the slope}), that the slope between these two vectors can be computed as $\textbf{slope}(e \sim y):= \frac{\Cov(e,y)}{\Var(y)}$. Thus, we are going to treat this definition as the corner-stone for deriving the choices of penalty term that can be included in the \textit{base model} (\ref{opt:base_model}), noting that $\Var(y)$ is constant in terms of the model choice, so the focus is only on the covariance term. To generalize the choices of covariance-based penalty terms, we denote them by $C(e,y)$. Hence, the regressivity-controlled learning problem can be written as   
\begin{align}
    \label{opt:general_penalized_model}
    f^*_{C}:=\arg\min_{f \in \H} \; \frac{1}{n} \sum_{i=1}^n \ell(f(x_i), y_i) + \mathcal{R}(f) + C(e, y),
\end{align}
where the exact form of $C(e,y)$ depends on the choice of constraints that want to be enforced in the model.

% \section{Soft-constrained Regression Models}
% Let $\D:=\{(x_i.y_i)\}_{i=1}^n$ be a given training dataset of $n$ properties and $p$ features associated with them (e.g., number of bedrooms, proximity to a bus stop, etc.), where $X\in \R^{n\times p}$ is the matrix of features and $y\in \R^n$ is the target vector (log-prices). Let $\H$ be a hypothesis class associated with some regression learning model (e.g. linear regression, gradient boosting, etc.), which contains functional mappings of the form $f:\X \rightarrow \Y$, where $\X$ is some feature space and $\Y$ is the target's space (log-prices space). In order to obtain the best fitting mapping $f^* \in \H$, we solve the complexity-regularized \textit{empirical risk minimization} problem 
% \begin{align}
%     \label{opt:base_model}
%      f^*:=\arg\min_{f \in \H} \; \frac{1}{n} \sum_{i=1}^n \ell(f(x_i), y_i) + \mathcal{R}(f) ,
% \end{align}
% where $\ell$ is some predefined loss function  (e.g., squared error), and $\mathcal{R}:\R^p \rightarrow \R$ is some complexity-control mapping (e.g., $\ell_1$-norm regularization), which is optional but commonly used in practice. In what follows, we are going refer to the solution of (\ref{opt:base_model}) as the \textit{base model}.

\medskip

\noindent\textbf{Covariance-penalized Regression.} In order to add regressivity-awareness to the base model (\ref{opt:base_model}), via linear trend correction between the Assessment-to-Sale ratios $r$ and the log-prices $y$, one of the natural choices of the penalty term $C(e,y)$ is to add a covariance regularization penalty. Specifically, since our target is a vector of log-prices $y$, we are going to penalize the empirical covariance between logarithm of the ratios (i.e., the residuals $e_i =   f(x_i) - y_i = \log r_i$) and the vector $y$, which is defined as $\widehat{\Cov}(e ,y)= \frac{1}{n}\sum_{i=1}^n (e_i - \Bar{e})(y_i - \Bar{y})$, where $\Bar{e}$ and $\Bar{y}$ refer to the average value of the residual vector and the log-price vector, respectively. We can enforce linear trend corrections in two ways with the covariance:
\begin{enumerate}
    \item \textbf{Squared Covariance:} Penalize any type of relationship (regressive or progressive) between $e$ and $y$, as
\begin{equation}
\everymath{\displaystyle}
\begin{array}{rcl}
 f_{\text{cov$^2$}}^*:=\arg\min_{f \in \H} \frac{1}{n} \sum_{i=1}^n \ell(f(x_i), y_i) + \frac{\rho}{2}\widehat{\text{Cov}}(e , y)^2 ,
\end{array}
\end{equation}

    \item \textbf{Negative Covariance:} Penalize only the negative trends between $e$ and $y$ (regressive patterns). 
\begin{equation}
\everymath{\displaystyle}
\begin{array}{rcl}
 f_{\text{neg-cov}}^*:=\arg\min_{f \in \H} \frac{1}{n} \sum_{i=1}^n \ell(f(x_i), y_i) - {\rho}\widehat{\text{Cov}}(e , y) ,
\end{array}
\end{equation}
\end{enumerate}
 In these two definitions, $\rho\geq 0$ referes to the penalization weight that dictates the soft-constraint level, and $\widehat{\text{Cov}}$ refers to the empirical covariance. 
 
 \textit{Note:} These formulations can be generalized to enforce any desired kind of linear independence between some transformation of the the predictions $f(x) \in \R^n$ and any other sensitive quantity $d\in \R^n$. 


\medskip

\noindent \textbf{A separable surrogate.} The covariance constraint has two main caveats: (i) it can hide positive-negative relationships that cancel out to get a close-to-zero correlation






































































(e.g., V-shaped trend), and (ii) if enforcing any kind of trend (regressive or progressive), it is not a separable on the different data-points (i.e., its coupling all the dataset), which is not a desirable property on some important ML settings (e.g., XGBoost, LightGBM, Neural Networks). To address these two issues we first note that we can make use of Jensen's inequality over the squared covariance term to get the upper bound $\widehat{\text{Cov}}(e, y)^2= \left(\frac{1}{n}\sum_{i=1}^{n}(e_i - \Bar{e})(y_i -\Bar{y}) \right)^2\leq  \frac{1}{n}\sum_{i=1}^{n}(e_i - \Bar{e})^2(y_i -\Bar{y})^2$. Further, note that this is just a weighted deviation correction of the residuals $e_i$ from their average $\Bar{e}$. Ideally, the average of the residuals is expected to lie near zero (i.e., $\Bar{e} \approx 0$), and the same would be expected from each residual $e_i$. Thus penalizing the residual deviations $(e_i - 0)^2$, with a focus on greater weights on extreme values, seems a natural choice of surrogate. Hence, we define the following surrogate function
\[
\text{Surr}(e,y) : = \frac{1}{n}\sum_{i=1}^n e_i^2(y_i - \Bar{y})^2,
\]
which instead of penalizing linear dependencies only, is also able to capture some nonlinear relationships between $e$ and $y$. Note that its also a separable term, allowing us to propose the datapoint-wise separable surrogate-penalized learning problem 




$f^*_{\text{cov}}$ in $\H$, we solve the general-form regression problem 

% $\widehat{\text{Cov}}(e, y)^2= \left(\frac{1}{n}\sum_{i=1}^{n}(e_i - \Bar{e})(y_i -\Bar{y}) \right)^2\leq  \frac{1}{n}\sum_{i=1}^{n}(e_i - \Bar{e})^2(y_i -\Bar{y})^2$. Further, since this upper bound is simply a weighted penalization over the deviations of $e$ from its average, and we want to enforce $\Bar{e}=0$, we can write the surrogate-regularized problem as   
\begin{equation}
\everymath{\displaystyle}
\begin{array}{rcl}
 f^*_{\text{surr}}:=\arg\min_{f \in \H} \frac{1}{n} \sum_{i=1}^n \ell(f(x_i), y_i) + \frac{\rho}{2}\text{Surr}(e , y) .
\end{array}
\end{equation}
We don't use this surrogate re-formulation of the learning problem just an alternative, but also as our main and suggested regularized model for usage in practice.

\noindent \textbf{Interpretation of the regularization terms.} We analyze the meaning of the regularization terms from a statistical point of view.

\begin{itemize}
    \item \textbf{Covariance Penalty:} Note
\end{itemize}


\subsection{Covariance-regularized Regression Models Examples}

\subsubsection{Linear regression}

Let $\Theta$ be the space of linear coefficient, i.e., the parameters $\theta \in \R^{p}$ that define linear mappings of the form $f(x_i) = \theta ^\top x_i, \forall_i \in [n]$. Commonly used examples of these linear mappings include linear regression, the LASSO, Ridge Regression, and Elastic Net. The loss function of these kind of linear models is usually the mean squared error (MSE) plus some complexity control over the parameter search space $\Theta$. We are going to analyze the closed-form solution for the two proposed regularization terms for the case of linear regression. Thus, the coefficients are the optimal solution of the learning problem
\[
\everymath{\displaystyle}
\begin{array}{rl}
  \theta_{\text{cov}}^*= &\arg\min_{\theta \in \Theta}  \frac{1}{2n} \norm{X\theta - y}_2^2 + \frac{\rho}{2} \left(\frac{1}{n}\sum_{i=1}^n (e_i - \Bar{e})(y_i - \Bar{y}) \right)^2 \\
  =& \left( \frac{1}{n} X^\top X + \frac{\rho}{n^2} (X^{\top}Hy)(X^{\top}Hy)^{\top}  \right)^{-1} \left( \frac{1}{n}X^\top y + \frac{\rho}{n^2}  (y^{\top}Hy) (X^{\top}Hy)  \right) , 
\end{array}
\]
where $H= I - \frac{1}{n} \textbf{e}\textbf{e}^\top$ is an auxiliary centering symmetric matrix, with $\textbf{e}\in \R^n$ being a vector of 1's. In the case of the surrogate regularization, the linear regression optimal coefficients are the closed-form-solution for the learning problem
\[
\everymath{\displaystyle}
\begin{array}{rl}
  \theta_{\text{surr}}^*= &\arg\min_{\theta \in \Theta}  \frac{1}{2n} \norm{X\theta - y}_2^2 + \frac{\rho}{2n} \sum_{i=1}^n e_i^2 (y_i - \Bar{y})^2  \\
  =& \left( \frac{1}{n} X^\top X + \frac{\rho}{n} X^{\top}DX  \right)^{-1} \left( \frac{1}{n}X^\top y + \frac{\rho}{n}  X^{\top}Dy  \right) , 
\end{array}
\]
where $D $ is a diagonal matrix such that  $D_{ii}= (y_i  - \Bar{y})^2, \forall i \in [n]$. In both of this cases, the linear regression predictive model is of the form $f(x_i)=x_i^\top \theta^*, \forall i \in [n]$. One can further extend the solution to regularized linear models like the LASSO and Ridge regression considering some complexity-control mapping $\mathcal{R}:\R^p \rightarrow \R$ (e.g., $\lambda\norm{\theta}_1$, with $\lambda>0$, for the LASSO), to be included in the objective. This would have a closed-form in the $\ell_2$-norm case, but not if we include the $\ell_1$-norm regularization.


\subsubsection{Second-Order Gradient Boosting Machines}

Given its state-of-the-art performance on real world problem (\todo{citation}), we are going to focus solely on providing the covariance-based modifications on second-order gradient boosting methods. Common and widely used (\todo{CITE}) algorithm examples of this family include XGBoost and LightGBM. We are not going to delve into every detail of each one of the implementations of this family, but on the general approach on how to implement the custom objective.
% \textbf{First-Order Gradient Boosting}

% \begin{algorithm}[!ht]
% \caption{{\color{red}(Sketch of)} Adversarial Gradient Boosting Algorithm}
% \label{alg:gradient_boosting}
% \begin{algorithmic}[1] % The '[1]' enables line numbering

%     % INPUTS
%     \Require Training set $\{(x_i, y_i)\}_{i=1}^N$, differentiable loss function $L(y, F(x))$, number of iterations $T$, learning rate $\eta$.

%     % INITIALIZATION
%     \State Initialize model with a constant value:
%     $ F_0(x) = \operatorname*{arg\,min}_{\gamma} \sum_{i=1}^N L(y_i, \gamma) $

%     % MAIN LOOP
%     \For{$t = 1$ to $T$}
        
%         \State \textbf{1. Compute pseudo-residuals:}
%         \For{$i = 1$ to $N$}
%             \State $ r_{it} = - \left[ \frac{\partial L(y_i, F(x_i))}{\partial F(x_i)} \right]_{F(x)=F_{t-1}(x)} \quad \text{ i.e. }(y_i - \hat{y}_{i})^2 \text{ in our case }$
%         \EndFor
%         \State \textbf{{\color{red}2. Compute some adversarial re-weighting vector $w_t\in \mathbb{R}^n$ (to be defined)}}
%         \State \textbf{3. Fit base learner:}
%         \State Fit a weak learner (e.g., regression tree) $h_t(x)$ to the re-weighted training data $\{(x_i, r_{it}, w_{it})\}_{i=1}^N$.
%         % \State \textbf{3. Compute step size:}
%         % \State Solve the 1D optimization problem to find optimal leaf weights:
%         % $ \gamma_m = \operatorname*{arg\,min}_{\gamma} \sum_{i=1}^N L(y_i, F_{m-1}(x_i) + \gamma h_m(x_i)) $
%         \State \textbf{4. Update model:}
%         $ F_t(x) = F_{t-1}(x) + \eta h_t(x) $
        
%     \EndFor

%     % OUTPUT
%     \State \Return Final model $F_T(x)$

% \end{algorithmic}
% \end{algorithm}

% \begin{algorithm}[!ht] % Use [t] or [htbp] depending on your paper template
% \caption{First-Order Gradient Boosting with Squared Covariance Penalty}
% \label{alg:cov_boosting}
% \begin{algorithmic}[1]
% \Require Training data $\mathcal{D} = \{(x_i,y_i)\}_{i=1}^n$, penalty $\rho \ge 0$, learning rate $\nu \in (0,1]$, rounds $M$
% \State \textbf{Initialize:} $\bar{y} \gets \frac{1}{n}\sum_{i=1}^n y_i$, \quad $f_0(x) \gets \bar{y}$
% \State Precompute centered targets: $c_i \gets y_i - \bar{y}$ \quad $\forall i \in \{1, \dots, n\}$
% \For{$m = 1, \dots, M$}
%     \State $f_i \gets f_{m-1}(x_i)$ \textbf{and} $e_i \gets y_i - f_i$ \quad $\forall i \in \{1, \dots, n\}$ \Comment{Predictions \& residuals}
%     \State $z \gets \frac{1}{n}\sum_{i=1}^n e_i c_i$ \Comment{Covariance statistic}
%     \State $\widetilde{r}_i \gets e_i + \rho z c_i$ \quad $\forall i \in \{1, \dots, n\}$ \Comment{Form pseudo-residuals}
%     \State Fit a base learner structure $T_m$ to $\{(x_i,\widetilde{r}_i)\}_{i=1}^n$
    
%     \If{$T_m$ is a generic weak learner with output vector $h \in \mathbb{R}^n$}
%         \State $q \gets \frac{1}{n}c^\top h$
%         \State $\gamma_m \gets \frac{\frac{1}{n}e^\top h + \rho z q}{\frac{1}{n}\|h\|_2^2 + \rho q^2}$ \Comment{Exact line-search step}
%         \State $f_m(\cdot) \gets f_{m-1}(\cdot) + \nu \gamma_m h(\cdot)$
%     \Else \Comment{Tree-based learner}
%         \State Let $I_1, \dots, I_J$ be the disjoint leaves of $T_m$
%         \For{$j = 1, \dots, J$}
%             \State $n_j \gets |I_j|, \quad b_j \gets \sum_{i \in I_j} e_i, \quad s_j \gets \sum_{i \in I_j} c_i, \quad q_j \gets s_j/n$
%         \EndFor
%         \State $D \gets \operatorname{diag}\!\left(\frac{n_1}{n}, \dots, \frac{n_J}{n}\right), \quad r \gets \frac{1}{n}b + \rho z q$
%         \State Solve the coupled leaf system: $(D + \rho qq^\top)w = r$
%         \State $f_m(\cdot) \gets f_{m-1}(\cdot) + \nu\,T_m(\cdot;w)$
%     \EndIf
% \EndFor
% \State \Return final model $f_M$
% \end{algorithmic}
% \end{algorithm}

% \textbf{Second-Order Gradient Boosting}

\begin{algorithm}[!ht]
\caption{First-order gradient boosting with squared covariance penalty}
\begin{algorithmic}[1]
\Require Training data $\D=\{(x_i,y_i)\}_{i=1}^n$, penalty weight $\rho \ge 0$, learning rate $\nu \in (0,1]$, number of rounds $M$
\State Compute $c_i \gets y_i - \bar y$, where $\bar y = \frac{1}{n}\sum_{i=1}^n y_i$
\State Initialize $f_0(x) \gets \bar y$
\For{$m = 1,\dots,M$}
    \State Compute current predictions $f_i \gets f_{m-1}(x_i)$ for all $i$
    \State Compute residuals $e_i \gets y_i - f_i$ for all $i$
    \State Compute covariance statistic
    \[
    z \gets \frac{1}{n}\sum_{i=1}^n e_i c_i
    \]
    \State Form pseudo-residuals
    \[
    \widetilde r_i \gets e_i + \rho z c_i
    \qquad \forall i=1,\dots,n
    \]
    \State Fit a base learner structure $T_m$ to $\{(x_i,\widetilde r_i)\}_{i=1}^n$
    \If{$T_m$ is a generic weak learner with output vector $h \in \R^n$}
        \State Compute
        \[
        q \gets \frac{1}{n}c^\top h
        \]
        \State Set exact line-search step
        \[
        \gamma_m \gets
        \frac{\frac{1}{n}e^\top h + \rho z q}
        {\frac{1}{n}\|h\|_2^2 + \rho q^2}
        \]
        \State Update
        \[
        f_m(\cdot) \gets f_{m-1}(\cdot) + \nu \gamma_m h(\cdot)
        \]
    \Else
        \State Let $I_1,\dots,I_J$ be the leaves of $T_m$
        \For{$j=1,\dots,J$}
            \State $n_j \gets |I_j|$
            \State $b_j \gets \sum_{i \in I_j} e_i$
            \State $s_j \gets \sum_{i \in I_j} c_i$
            \State $q_j \gets s_j/n$
        \EndFor
        \State Form
        \[
        D \gets \diag\!\left(\frac{n_1}{n},\dots,\frac{n_J}{n}\right),
        \qquad
        r \gets \frac{1}{n}b + \rho z q
        \]
        \State Solve the coupled leaf system
        \[
        (D+\rho qq^\top)w = r
        \]
        \State Update
        \[
        f_m(\cdot) \gets f_{m-1}(\cdot) + \nu\,T_m(\cdot;w)
        \]
    \EndIf
\EndFor
\State \Return $f_M$
\end{algorithmic}
\end{algorithm}

\begin{algorithm}[!h]
\caption{Exact Newton-style tree boosting with squared covariance penalty}
\begin{algorithmic}[1]
\Require Training data $\{(x_i,y_i)\}_{i=1}^n$, penalty $\rho \ge 0$, leaf regularization $\lambda \ge 0$, complexity penalty $\gamma \ge 0$, learning rate $\nu \in (0,1]$, number of rounds $M$
\State Compute $c_i \gets y_i - \bar y$ for all $i$
\State Initialize prediction function $f_0$
\For{$m = 1,\dots,M$}
    \State Compute current predictions $f_i \gets f_{m-1}(x_i)$
    \State Compute residuals $e_i \gets y_i - f_i$
    \State Compute
    \[
    z \gets \frac{1}{n}\sum_{i=1}^n e_i c_i
    \]
    \State Compute exact gradients
    \[
    g_i \gets \frac{f_i-y_i}{n} - \frac{\rho}{n} z c_i
    \qquad \forall i=1,\dots,n
    \]
    \State Search over candidate tree structures $T$
    \For{each candidate tree $T$ with leaves $I_1,\dots,I_J$}
        \For{$j=1,\dots,J$}
            \State $G_j \gets \sum_{i \in I_j} g_i$
            \State $n_j \gets |I_j|$
            \State $s_j \gets \sum_{i \in I_j} c_i$
        \EndFor
        \State Form
        \[
        D \gets \diag\!\left(\frac{n_1}{n}+\lambda,\dots,\frac{n_J}{n}+\lambda\right),
        \qquad
        \beta \gets \frac{\rho}{n^2},
        \qquad
        A_T \gets D + \beta ss^\top
        \]
        \State Compute exact tree gain
        \[
        \operatorname{Gain}(T)
        \gets
        \frac{1}{2}G^\top A_T^{-1}G - \gamma J
        \]
    \EndFor
    \State Select the best tree structure $T_m$
    \State For the selected tree, form $A_m = D + \beta ss^\top$
    \State Compute exact leaf weights
    \[
    w^\star \gets -A_m^{-1}G
    \]
    \State Update
    \[
    f_m(\cdot) \gets f_{m-1}(\cdot) + \nu\,T_m(\cdot;w^\star)
    \]
\EndFor
\State \Return $f_M$
\end{algorithmic}
\end{algorithm}





\subsection{In-processing Regression Problem}

\subsection{Post-processing Regression Problem}




Covariance-constrained models. 

% \section{Second-stage of approach? The robust model selection}



\section{Implementation}

\section{Multi-objective Model Selection}

% Pareto-optimal decisions

% (there exist more than one method to perform this model selection like the min-max approach, Pareto and beyond, etc.)

\textbf{Two-objective Robust Model Selection.} Let $\C := \{c_1, c_2, \ldots, c_m\}$ be a finite collection of $m > 1$ already trained models. Let $W:=\{w_1, w_2, \ldots, w_F\}$ be a finite collection of $F>1$ folds (or time windows), and $\{T_w, V_w\}_{w\in W}$ be a collection of training and validation sets obtained from the training set $\D$ by appropriately splitting it in its respective process (cross validation or rolling-origin). Further, let $\A\subset \R^{m \times W}$ be some accuracy results collection of OOS for each configuration $c \in \C$ on the validation set of each window $w \in W$. Let $\F \in \R^{m \times W}$ be the analogous set for a given fairness metric. Then, given an objective level of fairness $\tau >0$, one way of obtaining a robust Pareto-optimal solution is by sorting the results, secure the fairness level on every fold/windows $w \in W$, and select the best performing configuration in an aggregated sense. That is, solving the robust problem
\[
\everymath{\displaystyle}
\begin{array}{rcl}
c^* := \arg\min_{c \in \C} & \frac{1}{W}\sum_{i\in w}  \A_w(c) --- \ell(f(x_i) ,y_i;c) \\
\text{s.t.} & \F_{w}(c), & \forall w\in W .
\end{array}
\]



{ 
\begin{enumerate}
\item \textbf{Train:} $c\in \C$ configurations of the model, and $f_c $ the predictive model they produce. 
{\small
\begin{itemize}
        \item E.g.: Different models, hyperparameters, penalty $\rho>0$.
    \end{itemize}}
    \item  \textbf{Evaluate:} on $v\in\V$ validation folds, some 
    \begin{itemize}

    \item \underline{Prediction error} metric $\mathcal{E}_v(f_c)$ (e.g., MSE), and some
    
    \item \underline{Regressivity} metric ${\F_v}(f_c)$ (e.g., PRD).
    \end{itemize}
        \item \textbf{Calibration Alternative:} Solve the \underline{sorting problem} (2-objectives example)
\[
\everymath{\displaystyle}
\begin{array}{rclccc}
(MS_\tau) \quad f_c^*:=\arg\min_{f_c : c \in \C} & \frac{1}{|\V|} \sum_{v\in \V} \mathcal{E}_v(f_c) & \\
\text{s.t.} &   \F_v(f_c) \leq \tau, &\forall v \in \V,
\end{array}
\]
or some other aggregation of the results $\F_v(f_c)$, for $v \in \V$. 
    \end{enumerate}
% \medskip
% \vspace{-2mm}
}
Note that the accuracy and fairness measures used in this model, don't even need to be the loss and covariance metrics of the originally-trained model, but any accuracy, fairness, or adjacent metric that one may desire to work with. 


\subsection{Convex Stacking of Models}

We can further extend $(MS_\tau)$ to get a Pareto-optimal stacking of models. That is, include possible combinations of the prediction vectors $f_c$ as feasible solutions. That is, solving
\[
\everymath{\displaystyle}
\begin{array}{rclccc}
(\text{Stack-}MS_\tau) \quad f_{w^*}:=\arg\min_{w \in \R^{|\C|}} & \frac{1}{|\V|} \sum_{v\in \V} \mathcal{E}_v(f_{w}) & \\
\text{s.t.} &   \F_v(f_w) \leq \tau, &\forall v \in \V,\\
& \sum_{c\in \C} w_c = 1, \\
& w \geq 0. 
\end{array}
\]
where $f_w = \sum_{c\in \C} w_cf_c $ is a stack of models/configurations. Note that, if $\mathcal{E}$ and $\F$ are convex mappings of $f$ (fixed model), then the model stacking metric preserves the convexity. Further, in practice, many of the regressivity metrics constraints can be re-written in a form that preserves the convexity of the problem, in terms of the weights $w \in \R_+^{|\C|}$; i.e, they allow a convex formulation of (Stack-$SP_\tau)$.




% Robust Pareto-based model selection. 

\subsection{Formulation Examples of Performance Metrics}

\subsubsection{General Metrics Formulations}

\begin{itemize}
    \item \textbf{Median Ratio:} The median of the Assessment-to-Sales ratio is one of the general metrics of interest defined by IAAO. Most of the more specific measures depend on it.
    $$
    \tilde r := \text{median}(r) \approx 1 
    $$

    We need to introduce an auxiliary variable, $r_m$, that takes the role of the median ratio through the following constraints
    \[
    \everymath{\displaystyle}
    \begin{array}{cl}
       r_i + (1-z_i) M \leq r_m \leq r_i + z_i M , & \forall i \in [n], \\ 
       \sum_{i=1}^n z_i = \frac{n}{2} + \mathds{1}_{\left\{ \frac{n }{ 2} \;\in \; \mathbb{Z} \right\}}, 
    \end{array}
    \]
and then we actually constraint the value of this median value as 
\[
  1 - \epsilon \leq r_m \leq 1 + \epsilon, 
\]
for some tolerance of deviation $\epsilon$ from 1.
\end{itemize}

\subsubsection{Regressivity Metrics Constraint Formulations}

\begin{itemize}
    \item \textbf{Price-Related Differential (PRD):}  The most relevant metric for measuring regressivity, at least before 2025, was the Price-Related Differential (PRD). Lets recall that, given a predictive model $f\in \H$, the PRD is defined as the ratio between the \todo{Confirm name} Assessment-to-Price  ratio's average and their weighted average, i.e., 
    $$PRD(f) := \frac{\frac{1}{n} \sum_{i=1}^n r_i}{\sum_{i=1}^n f(x_i) / \sum_{i=1}^n y_i} = \frac{\frac{1}{n} \sum_{i=1}^n f(x_i)/y_i}{\sum_{i=1}^n f(x_i) / \sum_{i=1}^n y_i}. $$ For  If we add this term and its reference interval as a constraint in $(MS_\tau)$, then we simply filter by those $c\in \C$ that satisfy $PRD(f_c) \in [0.98, 1.03]$. On the other hand, for (Stack-$MS_\tau$), we can re-formulate the constraint (given $f(x_i), y_i > 0$) as  
    \vspace{-2mm}
    \begin{align*}
     0.98 \leq PRD(f_w) \leq 1.03   \iff & 0.98\frac{  \sum_{i=1}^n f_w(x_i)}{ \sum_{i=1}^n y_i}    \leq \frac{1}{n}\sum_{i=1}^n \frac{f_w(x_i)}{ y_i} \leq 1.03 \frac{ \sum_{i=1}^n f_w(x_i)}{ \sum_{i=1}^n y_i},   
    \end{align*}
    which is linear (convex) over the model weights $w \in \R_{+}^{|\C|}$.

    \item \textbf{Price-Related Bias (PRB).} 

    \begin{align*}
       PRB:= &\text{slope}\left( \frac{r_i - \tilde r }{\tilde r} \sim \log_2\left( \frac{1}{2} ( MV_i + AV_i/\tilde r ) \right)  \right)\\
     =& \frac{\sum_{i=1}^n (\log_2\left( \frac{1}{2} ( MV_i + AV_i/\tilde r ) \right) - \Bar{x})(y_i - \Bar{y})}{\sum_{i=1}^n (x_i - \Bar{x})^2}
    \end{align*}
    
\end{itemize}









\section{(Optional) important points of the approach. Other have empirical examination or case of study can be an option.}






\section{Experiments}

% \section{Extensions}

\section{Discussion/Conclusions, Limitations, Extensions \& Future Work}

\subsection{Future Work}

Same direction: control the relationship, but now aware of the sub-market clusters, which can create nonlinear relationships when asking for too much overall linear correction (see image of 0-correlation plot of surrogate penalty).

Interpretable \& fair models. Apply this to multiple cases of assessors offices, neural networks as n alternative methodology because of their predictive power (?); in particular, their ability to handle Geospatial data (e.g., apply an embedding using the information or methodology provided by AlphaEarth of Google). Another alternative is to test the foundational models (e.g., tabular models). Try to incorpore it into Random Forest/Bagging, which is also widely used apparently.

\section{MY NOTES}

Paper structure references. INFORMS seems to use a lot of plots and tables in general (make sense since they care about the analytics part).
\begin{itemize}
    \item INFORMS data science: https://pubsonline.informs.org/doi/epdf/10.1287/ijds.2024.0029
    \item INFORMS managements science: https://pubsonline.informs.org/doi/epdf/10.1287/mnsc.2022.00345
    \item INFORM journal on optimization: https://arxiv.org/pdf/2412.17116
\end{itemize}



\noindent Haihao whiteboard:
\begin{itemize}
    \item Introduction: motivation, contribution, related work.
    \item Setting: Examples, assumptions, etc.
    \item Methodology
    \item Experiments 
    \item Extension (2-pages at most)
    \item Conclusions
\end{itemize}

\noindent Swati Comments / Suggestions:

\begin{itemize}
    \item What is novel in this research (theory and applications-wise)? Has anything like this be done already?
    \item Good thing: it looked like I owned the project. Lets keep it that way for the GE.
    \item It has to be a journal-level paper. Good venues: Journal of Machine Learning Research (JMLR) or Management Sciences. 
    \item Paper Critique: adjacent to what I've done. Maybe ML, fairness, etc. The idea is to be able to read, summarize, critique (Not just chatGPT-style). Even go on details, proofs, technical stuff. I can do extra research on the topic of the paper.
    \item Important: GE as a message of what I've been doing during the phd.
    \item Presentation (oral-exam): Be prepared for presenting + questions at the same time. Bad things about past presentations: too much notation, not be able to answer extra questions, etc.
    \item Benchmark v/s more complex correlation/covariance methods. They don't work/generalize as good. Show it.
    \item Extras: copy of the paper \& slides in physical for reading. 
\end{itemize}

\printbibliography %Prints bibliography


% APPENDIX
\newpage

\appendix

\section{Algorithmic Details for Covariance-Penalized Boosting}
\label{app:boosting_algorithms}

This appendix collects pseudocode for boosted-tree implementations of the squared covariance penalty introduced in Section~\ref{sec:covariance_penalizations}. Throughout, we use the manuscript's sign convention
\[
e_i = f(x_i)-y_i,
\qquad
c_i = y_i-\bar y,
\qquad
z = \widehat{\Cov}_n(e,y)=\frac{1}{n}\sum_{i=1}^n e_i c_i.
\]
Under this convention, the squared covariance penalty is
\[
\Psi_n^{\mathrm{cov}^2}(f)=\frac{1}{2}z^2.
\]

\subsection{Exact Newton-style tree boosting with squared covariance penalty}
\label{app:exact_newton_cov_boosting}

The following algorithm implements the exact leaf-wise Newton update discussed in Section~\ref{subsec:tree_boosting}. This is the most relevant pseudocode to retain if only one boosted-tree algorithm is included in the appendix.

\begin{algorithm}[!ht]
\caption{Exact Newton-style tree boosting with squared covariance penalty}
\label{alg:exact_newton_cov_boosting}
\begin{algorithmic}[1]
\Require Training data $\D=\{(x_i,y_i)\}_{i=1}^n$, penalty weight $\rho \ge 0$, leaf regularization $\lambda \ge 0$, tree complexity penalty $\gamma \ge 0$, learning rate $\nu \in (0,1]$, number of rounds $M$
\State Compute $\bar y \gets \frac{1}{n}\sum_{i=1}^n y_i$
\State Compute $c_i \gets y_i-\bar y$ for all $i=1,\dots,n$
\State Initialize $f_0(x) \gets \bar y$
\For{$m=1,\dots,M$}
    \State Compute current predictions $f_i \gets f_{m-1}(x_i)$ for all $i=1,\dots,n$
    \State Compute residuals $e_i \gets f_i-y_i$ for all $i=1,\dots,n$
    \State Compute covariance statistic
    \[
    z \gets \frac{1}{n}\sum_{i=1}^n e_i c_i
    \]
    \State Compute exact gradients
    \[
    g_i \gets \frac{e_i}{n} + \frac{\rho}{n} z c_i
    \qquad \forall i=1,\dots,n
    \]
    \State Search over candidate tree structures $T$
    \For{each candidate tree $T$ with leaves $I_1,\dots,I_J$}
        \For{$j=1,\dots,J$}
            \State $G_j \gets \sum_{i\in I_j} g_i$
            \State $n_j \gets |I_j|$
            \State $s_j \gets \sum_{i\in I_j} c_i$
        \EndFor
        \State Form
        \[
        D \gets \diag\!\left(\frac{n_1}{n}+\lambda,\dots,\frac{n_J}{n}+\lambda\right),
        \qquad
        \beta \gets \frac{\rho}{n^2},
        \qquad
        A_T \gets D + \beta ss^\top
        \]
        \State Compute exact tree gain
        \[
        \operatorname{Gain}(T)
        \gets
        \frac{1}{2}G^\top A_T^{-1}G - \gamma J
        \]
    \EndFor
    \State Select the best tree structure $T_m$
    \State For the selected tree, form $A_m \gets D + \beta ss^\top$
    \State Compute optimal leaf weights
    \[
    w^\star \gets -A_m^{-1}G
    \]
    \State Update
    \[
    f_m(\cdot) \gets f_{m-1}(\cdot) + \nu\,T_m(\cdot;w^\star)
    \]
\EndFor
\State \Return final predictor $f_M$
\end{algorithmic}
\end{algorithm}

\subsection{Optional first-order boosting variant}
\label{app:first_order_cov_boosting}

For completeness, we also record a first-order boosting variant based on fitting base learners to the negative gradient of the penalized objective. This algorithm is mainly pedagogical, and can be omitted if a shorter appendix is preferred.

\begin{algorithm}[!ht]
\caption{First-order gradient boosting with squared covariance penalty}
\label{alg:first_order_cov_boosting}
\begin{algorithmic}[1]
\Require Training data $\D=\{(x_i,y_i)\}_{i=1}^n$, penalty weight $\rho \ge 0$, learning rate $\nu \in (0,1]$, number of rounds $M$
\State Compute $\bar y \gets \frac{1}{n}\sum_{i=1}^n y_i$
\State Compute $c_i \gets y_i-\bar y$ for all $i=1,\dots,n$
\State Initialize $f_0(x) \gets \bar y$
\For{$m=1,\dots,M$}
    \State Compute current predictions $f_i \gets f_{m-1}(x_i)$ for all $i=1,\dots,n$
    \State Compute residuals $e_i \gets f_i-y_i$ for all $i=1,\dots,n$
    \State Compute covariance statistic
    \[
    z \gets \frac{1}{n}\sum_{i=1}^n e_i c_i
    \]
    \State Form negative-gradient targets
    \[
    u_i \gets -\bigl(e_i + \rho z c_i\bigr)
    \qquad \forall i=1,\dots,n
    \]
    \State Fit a base learner structure $T_m$ to $\{(x_i,u_i)\}_{i=1}^n$
    \If{$T_m$ is a generic weak learner with output vector $h \in \R^n$}
        \State Compute
        \[
        q \gets \frac{1}{n}c^\top h
        \]
        \State Set exact line-search step
        \[
        \gamma_m
        \gets
        -
        \frac{\frac{1}{n}e^\top h + \rho z q}
        {\frac{1}{n}\|h\|_2^2 + \rho q^2}
        \]
        \State Update
        \[
        f_m(\cdot) \gets f_{m-1}(\cdot) + \nu\,\gamma_m h(\cdot)
        \]
    \Else
        \State Let $I_1,\dots,I_J$ be the leaves of $T_m$
        \For{$j=1,\dots,J$}
            \State $n_j \gets |I_j|$
            \State $b_j \gets \sum_{i\in I_j} e_i$
            \State $s_j \gets \sum_{i\in I_j} c_i$
            \State $q_j \gets s_j/n$
        \EndFor
        \State Form
        \[
        D \gets \diag\!\left(\frac{n_1}{n},\dots,\frac{n_J}{n}\right),
        \qquad
        q \gets (q_1,\dots,q_J)^\top,
        \qquad
        \eta \gets \frac{1}{n}b + \rho z q
        \]
        \State Solve the coupled leaf system
        \[
        (D+\rho qq^\top)w = -\eta
        \]
        \State Update
        \[
        f_m(\cdot) \gets f_{m-1}(\cdot) + \nu\,T_m(\cdot;w)
        \]
    \EndIf
\EndFor
\State \Return final predictor $f_M$
\end{algorithmic}
\end{algorithm}


% \newpage
% \section{Bulletpoints}

% \begin{itemize}
%     \item \textbf{Why this problem?}
%     \begin{enumerate}
%         \item Mass Appraisal with ML has been recently advancing. Examples are the CCAO, 
%         \item Mistakes on it can generate highly unfair situations in terms of taxation; specially if regressive: low-valued houses are overcharged and high-values houses are undercharged (e.g.
%     \end{enumerate}

    
%     \item \textbf{Why this approach?}
% \end{itemize}

% ==========================================
% Table 1: Validation Split
% ==========================================
\begin{table}[!htbp]
\centering
\renewcommand{\arraystretch}{1.5} % Extra stretch for stacked cell contents
\caption{Validation Results: Metrics in bold show improvement over LGBMRegressor. Colored percentages in parentheses under other models indicate percentual change (improvement / worsening) compared to LGBMRegressor.}
\label{tab:validation_metrics}
\resizebox{\textwidth}{!}{%
\begin{tabular}{ll|rr|rr|rrr}
\toprule
Model & Config ID & R2 & MAE & COD & COV (IAAO) & VEI & PRD & PRB \\
\midrule
LinearRegression & fd63507d & 0.668 & \$92,314 & 32.716 & 0.490 & 8.425\% & 1.026 & 0.059 \\
LGBMRegressor & 0c59889c & 0.889 & \$58,430 & 23.283 & 0.377 & -19.961\% & 1.066 & -0.069 \\
\midrule
SINGLE\_OPT\_REQ\_ROBUST & - & 
\begin{tabular}{@{}c@{}} 0.881 \\ \small\color{worsRed}{(-0.90\%)} \end{tabular} & 
\begin{tabular}{@{}c@{}} \$63,918 \\ \small\color{worsRed}{(-9.39\%)} \end{tabular} & 
\begin{tabular}{@{}c@{}} 24.465 \\ \small\color{worsRed}{(-5.08\%)} \end{tabular} & 
\begin{tabular}{@{}c@{}} \textbf{0.361} \\ \small\color{impGreen}{(+4.24\%)} \end{tabular} & 
\begin{tabular}{@{}c@{}} \textbf{2.538\%} \\ \small\color{impGreen}{(+87.29\%)} \end{tabular} & 
\begin{tabular}{@{}c@{}} \textbf{1.009} \\ \small\color{impGreen}{(+86.36\%)} \end{tabular} & 
\begin{tabular}{@{}c@{}} \textbf{0.026} \\ \small\color{impGreen}{(+62.32\%)} \end{tabular} \\
SINGLE\_OPT\_UTOPIA\_ROBUST & - & 
\begin{tabular}{@{}c@{}} 0.881 \\ \small\color{worsRed}{(-0.90\%)} \end{tabular} & 
\begin{tabular}{@{}c@{}} \$63,918 \\ \small\color{worsRed}{(-9.39\%)} \end{tabular} & 
\begin{tabular}{@{}c@{}} 24.465 \\ \small\color{worsRed}{(-5.08\%)} \end{tabular} & 
\begin{tabular}{@{}c@{}} \textbf{0.361} \\ \small\color{impGreen}{(+4.24\%)} \end{tabular} & 
\begin{tabular}{@{}c@{}} \textbf{2.538\%} \\ \small\color{impGreen}{(+87.29\%)} \end{tabular} & 
\begin{tabular}{@{}c@{}} \textbf{1.009} \\ \small\color{impGreen}{(+86.36\%)} \end{tabular} & 
\begin{tabular}{@{}c@{}} \textbf{0.026} \\ \small\color{impGreen}{(+62.32\%)} \end{tabular} \\
\bottomrule
\end{tabular}%
}
\end{table}

\vspace{3em}

% ==========================================
% Table 2: Test Split
% ==========================================
\begin{table}[!htbp]
\centering
\renewcommand{\arraystretch}{1.5} % Extra stretch for stacked cell contents
\caption{Test Results: Metrics in bold show improvement over LGBMRegressor. Colored percentages in parentheses under other models indicate percentual change (improvement / worsening) compared to LGBMRegressor.}
\label{tab:test_metrics}
\resizebox{\textwidth}{!}{%
\begin{tabular}{ll|rr|rr|rrr}
\toprule
Model & Config ID & R2 & MAE & COD & COV (IAAO) & VEI & PRD & PRB \\
\midrule
LinearRegression & fd63507d & 0.762 & \$86,244 & 25.013 & 0.369 & -13.527\% & 1.039 & -0.011 \\
LGBMRegressor & 0c59889c & 0.882 & \$72,672 & 23.437 & 0.359 & -37.257\% & 1.088 & -0.120 \\
\midrule
STACKING\_OPT\_AVG\_ROBUST & - & 
\begin{tabular}{@{}c@{}} \textbf{0.883} \\ \small\color{impGreen}{(+0.11\%)} \end{tabular} & 
\begin{tabular}{@{}c@{}} \$73,080 \\ \small\color{worsRed}{(-0.56\%)} \end{tabular} & 
\begin{tabular}{@{}c@{}} \textbf{22.823} \\ \small\color{impGreen}{(+2.62\%)} \end{tabular} & 
\begin{tabular}{@{}c@{}} \textbf{0.335} \\ \small\color{impGreen}{(+6.69\%)} \end{tabular} & 
\begin{tabular}{@{}c@{}} \textbf{-13.100\%} \\ \small\color{impGreen}{(+64.84\%)} \end{tabular} & 
\begin{tabular}{@{}c@{}} \textbf{1.040} \\ \small\color{impGreen}{(+54.55\%)} \end{tabular} & 
\begin{tabular}{@{}c@{}} \textbf{-0.020} \\ \small\color{impGreen}{(+83.33\%)} \end{tabular} \\
STACKING\_OPT\_UTOPIA\_ROBUST & - & 
\begin{tabular}{@{}c@{}} 0.879 \\ \small\color{worsRed}{(-0.34\%)} \end{tabular} & 
\begin{tabular}{@{}c@{}} \$76,258 \\ \small\color{worsRed}{(-4.94\%)} \end{tabular} & 
\begin{tabular}{@{}c@{}} 23.636 \\ \small\color{worsRed}{(-0.85\%)} \end{tabular} & 
\begin{tabular}{@{}c@{}} \textbf{0.345} \\ \small\color{impGreen}{(+3.90\%)} \end{tabular} & 
\begin{tabular}{@{}c@{}} \textbf{-3.157\%} \\ \small\color{impGreen}{(+91.53\%)} \end{tabular} & 
\begin{tabular}{@{}c@{}} \textbf{1.027} \\ \small\color{impGreen}{(+69.32\%)} \end{tabular} & 
\begin{tabular}{@{}c@{}} \textbf{0.014} \\ \small\color{impGreen}{(+88.33\%)} \end{tabular} \\
SINGLE\_OPT\_REQ\_ROBUST & - & 
\begin{tabular}{@{}c@{}} 0.878 \\ \small\color{worsRed}{(-0.45\%)} \end{tabular} & 
\begin{tabular}{@{}c@{}} \$77,163 \\ \small\color{worsRed}{(-6.18\%)} \end{tabular} & 
\begin{tabular}{@{}c@{}} 23.929 \\ \small\color{worsRed}{(-2.10\%)} \end{tabular} & 
\begin{tabular}{@{}c@{}} \textbf{0.348} \\ \small\color{impGreen}{(+3.06\%)} \end{tabular} & 
\begin{tabular}{@{}c@{}} \textbf{-2.237\%} \\ \small\color{impGreen}{(+94.00\%)} \end{tabular} & 
\begin{tabular}{@{}c@{}} \textbf{1.025} \\ \small\color{impGreen}{(+71.59\%)} \end{tabular} & 
\begin{tabular}{@{}c@{}} \textbf{0.019} \\ \small\color{impGreen}{(+84.17\%)} \end{tabular} \\
SINGLE\_OPT\_UTOPIA\_ROBUST & - & 
\begin{tabular}{@{}c@{}} 0.878 \\ \small\color{worsRed}{(-0.45\%)} \end{tabular} & 
\begin{tabular}{@{}c@{}} \$77,163 \\ \small\color{worsRed}{(-6.18\%)} \end{tabular} & 
\begin{tabular}{@{}c@{}} 23.929 \\ \small\color{worsRed}{(-2.10\%)} \end{tabular} & 
\begin{tabular}{@{}c@{}} \textbf{0.348} \\ \small\color{impGreen}{(+3.06\%)} \end{tabular} & 
\begin{tabular}{@{}c@{}} \textbf{-2.237\%} \\ \small\color{impGreen}{(+94.00\%)} \end{tabular} & 
\begin{tabular}{@{}c@{}} \textbf{1.025} \\ \small\color{impGreen}{(+71.59\%)} \end{tabular} & 
\begin{tabular}{@{}c@{}} \textbf{0.019} \\ \small\color{impGreen}{(+84.17\%)} \end{tabular} \\
\bottomrule
\end{tabular}%
}
\end{table}




\end{document}



\end{document}
